% !TeX root = ../PIFB2.tex
\clearpage
\section{Proofs for the Collective Evidence Identity, Latent Lie-Algebra Links and Coarse Link Laws}\label{app:latent_link_proofs}

This appendix proves the statements of Sections~\ref{sec:free_energy}, \ref{sec:latent_lie_links} and~\ref{sec:rg_status} that concern the collective evidence identity, latent Lie-algebra links and coarse link laws. It uses the following standard facts with references: the chain rule for relative entropy~\cite{CoverThomas2006} and the disintegration of laws on standard Borel spaces~\cite{Kallenberg2021}, the elementary theory of exponential families~\cite{Brown1986} and of convex conjugates~\cite{Rockafellar1970}, the naturality of the exponential map and the Baker--Campbell--Hausdorff series~\cite{Hall2015}, the matrix exponential and logarithm~\cite{Higham2008}, and Culver's criterion for real matrix logarithms~\cite{culver1966existence}.

\paragraph{Conventions.}
All laws live on standard Borel spaces. For laws $Q$ and $P$, $D_{\mathrm{KL}}(Q\|P)=\int\log(dQ/dP)dQ$ when $Q\ll P$, and $D_{\mathrm{KL}}(Q\|P)=+\infty$ otherwise. In the first case the negative part of the integrand has $Q$-integral at most $1/e$, because $\int_{\{dQ/dP<1\}}-\log(dQ/dP)dQ=\int_{\{dQ/dP<1\}}-(dQ/dP)\log(dQ/dP)dP$ and $-u\log u\leq1/e$ on $[0,1]$; hence $D_{\mathrm{KL}}(Q\|P)\in[0,+\infty]$, with value zero exactly when $Q=P$. If $Q$ and $P$ have densities $q$ and $p$ with respect to a $\sigma$-finite reference, then $Q\ll P$ exactly when $Q(\{p=0\})=0$, since a $P$-null set $A$ has $A\cap\{p>0\}$ null for the reference. The chain rule states that for laws $Q$ and $P$ on a product $\mathsf Y_1\times\mathsf Y_2$ with first marginals $Q_1,P_1$ and conditional kernels $Q_{2|1},P_{2|1}$, the conditional divergence $y_1\mapsto D_{\mathrm{KL}}(Q_{2|1}(y_1,\cdot)\|P_{2|1}(y_1,\cdot))$ is measurable and
\begin{equation}
D_{\mathrm{KL}}(Q\|P)=D_{\mathrm{KL}}(Q_1\|P_1)+\int D_{\mathrm{KL}}\bigl(Q_{2|1}(y_1,\cdot)\big\|P_{2|1}(y_1,\cdot)\bigr)Q_1(dy_1)
\quad\text{in }[0,+\infty].
\label{eq:llp_chain_rule}
\end{equation}
Matter coordinates, including source labels when present, are written $y$. For a link family $\Theta$ the held-link population law is $\mathbb P_\Theta(dy,do)=P^Y_\Theta(dy)K_\Theta(y,do)$, with the proper matter prior $P^Y_\Theta$, which is $P_0\Pi$ of Eq.~\eqref{eq:population_joint} at the links $\Theta$, and a normalized record kernel with density $\ell_o(y;\Theta)$ against a reference $\nu^O$ on records, jointly measurable in $(\Theta,y,o)$. At a record $o$ with $0<p(o\mid\Theta)<\infty$ its evidence, posterior and free energy are
\begin{equation}
\begin{gathered}
p(o\mid\Theta)=\int\ell_o(y;\Theta)P^Y_\Theta(dy),\qquad
\Pi_\Theta=\frac{\ell_o(\cdot;\Theta)P^Y_\Theta}{p(o\mid\Theta)},\\
\mathcal F_V[Q';X,\Theta,o]=D_{\mathrm{KL}}(Q'\|\Pi_\Theta)-\log p(o\mid\Theta).
\end{gathered}
\label{eq:llp_held_objects}
\end{equation}
At $\Theta(\mathring\omega)$ we write $p(o\mid\mathring\omega)$, $\Pi_{\mathring\omega}$ and $P^Y_{\mathring\omega}$. On $\mathcal T=\mathfrak g^{E_+}$, gradients, Hessians, adjoints, traces and determinants refer to the product of the inner products chosen on its summands, its norm $\|\cdot\|$ and the Lebesgue measure normalized by it; on matrices $\|\cdot\|$ is the operator norm unless a subscript says otherwise. Constants $c,c_1,c_2,\ldots$ are local to one statement.

\subsection{The collective evidence identity}\label{app:llp_evidence}

Let $P^Y$ be a proper prior on the matter space, $\ell_o\geq0$ a measurable record likelihood with evidence $Z_o=\int\ell_odP^Y\in(0,\infty)$, and $\Pi_o=Z_o^{-1}\ell_oP^Y$ the posterior. For every law $Q$ on the matter space put
\begin{equation}
\mathcal F_V[Q;X,o]=D_{\mathrm{KL}}(Q\|\Pi_o)-\log Z_o\in[-\log Z_o,+\infty].
\label{eq:llp_extended_vfe}
\end{equation}
For Eq.~\eqref{eq:collective_vfe}, $P^Y=P_0\Pi$ and $\Pi_o=\mathbb P(\cdot\mid o,X)$.

\begin{proposition}[Collective evidence identity]\label{prop:llp_evidence}
(a) $\mathcal F_V[Q;X,o]\geq-\log Z_o$, with equality exactly when $Q=\Pi_o$.
(b) If $D_{\mathrm{KL}}(Q\|P^Y)<\infty$ and $\mathbb E_Q\max(\log\ell_o,0)<\infty$, then $\mathcal F_V[Q;X,o]=D_{\mathrm{KL}}(Q\|P^Y)-\mathbb E_Q\log\ell_o$ in $(-\infty,+\infty]$, where $\log0=-\infty$ and $\mathbb E_Q\log\ell_o\in[-\infty,+\infty)$.
(c) If $Q$ and $P^Y$ have densities $q$ and $p^Y$ with respect to a $\sigma$-finite reference, and $\mathbb E_Q|\log q|$ and $\mathbb E_Q|\log(\ell_op^Y)|$ are finite, then $\mathcal F_V[Q;X,o]=\mathbb E_Q[\log q-\log(\ell_op^Y)]$.
\end{proposition}

\begin{proof}
Part (a) is the nonnegativity of relative entropy and its zero case. For (b), let $M=\ell_oP^Y$, a measure of total mass $Z_o$ with $\Pi_o=M/Z_o$. Suppose first that $Q\ll M$, and let $r=dQ/dM$. Then $dQ/d\Pi_o=Z_or$, and $\int_{\{r<1\}}-\log r dQ=\int_{\{r<1\}}-r\log r dM\leq Z_o/e$, so $D_{\mathrm{KL}}(Q\|\Pi_o)=\log Z_o+\int\log r dQ$ with an integral in $(-\infty,+\infty]$. Since $M\ll P^Y$, $dQ/dP^Y=r\ell_o$, and $Q(\{\ell_o=0\})=\int_{\{\ell_o=0\}}r dM=0$, so $\log r=\log(dQ/dP^Y)-\log\ell_o$ holds $Q$-almost surely. The first term is $Q$-integrable because $D_{\mathrm{KL}}(Q\|P^Y)$ is finite and its negative part is integrable, and the second has a $Q$-integrable negative part by hypothesis, so the integral splits and gives the claim. Suppose next that $Q\not\ll M$, and choose $S$ with $M(S)=0<Q(S)$. Then $\Pi_o(S)=0$, so the left side is $+\infty$. Since $\int_S\ell_odP^Y=0$, $\ell_o$ vanishes $P^Y$-almost everywhere on $S$, hence $Q$-almost everywhere on $S$ because $Q\ll P^Y$; thus $\mathbb E_Q\log\ell_o=-\infty$ and the right side is $+\infty$ as well. For (c), $\Pi_o$ has the density $\ell_op^Y/Z_o$, which is positive $Q$-almost surely because $\log(\ell_op^Y)$ is $Q$-integrable, so $Q\ll\Pi_o$ with $dQ/d\Pi_o=Z_oq/(\ell_op^Y)$ $Q$-almost surely. Its logarithm is a sum of three $Q$-integrable terms, and integration gives $D_{\mathrm{KL}}(Q\|\Pi_o)=\log Z_o+\mathbb E_Q[\log q-\log(\ell_op^Y)]$.
\end{proof}

\subsection{The loop defect and the generator prior}\label{app:llp_prior}

Let $\rho:G\to\mathrm{GL}(\mathbb V)$ be a smooth representation on $\mathbb V=\mathbb R^{n_V}$ with differential $\mathrm d\rho$, and write $\mathsf G=\mathsf G^{\mathrm{fib}}$ for the positive-definite fiber matrix of Eq.~\eqref{eq:pifb_metric_loop_defect}. For matrices on $\mathbb V$ put $\langle\mathbf H,\mathbf H'\rangle_{\mathsf G}=\operatorname{tr}(\mathsf G^{-1}\mathbf H^\top\mathsf G\mathbf H')$ and $\mathbf H^*=\mathsf G^{-1}\mathbf H^\top\mathsf G$, so that $\langle\mathbf H,\mathbf H'\rangle_{\mathsf G}=\operatorname{tr}(\mathbf H^*\mathbf H')$; a different subscript on the norm names the matrix used in place of $\mathsf G$. Two identities of Lie theory are used repeatedly~\cite{Hall2015}: $\rho(\exp\xi)=\exp(\mathrm d\rho\xi)$, and $h\exp(\xi)h^{-1}=\exp(\operatorname{Ad}_h\xi)$ with $\mathrm d\rho(\operatorname{Ad}_h\xi)=\rho(h)\mathrm d\rho(\xi)\rho(h)^{-1}$. For a closed walk $p=(e_0,\ldots,e_{r-1})$ the reference holonomy $\mathring\Theta_p$ is the ordered product of Eq.~\eqref{eq:pifb_metric_loop_defect} with the reversal convention $\mathring\omega_{\bar e}=-\mathring\omega_e$, and the cycle map is $\mathbf C_p\mathring\omega=\sum_{a=0}^{r-1}\mathring\omega_{e_a}$, which counts each traversal with its orientation.

\begin{lemma}[Trace form]\label{lem:llp_trace_form}
(a) $\langle\cdot,\cdot\rangle_{\mathsf G}$ is an inner product with $\|\mathbf H\mathbf H'\|_{\mathsf G}\leq\|\mathbf H\|_{\mathsf G}\|\mathbf H'\|_{\mathsf G}$ and $\|I\|_{\mathsf G}=n_V^{1/2}$, and left or right multiplication by $\mathbf H$ has as adjoint the left or right multiplication by $\mathbf H^*$.
(b) For positive-definite $\mathsf G'$ and $\mathsf G''$, the norms $\|\cdot\|_{\mathsf G'}$ and $\|\cdot\|_{\mathsf G''}$ agree on every rank-one matrix exactly when $\mathsf G''$ is a positive multiple of $\mathsf G'$, and then they agree on all matrices.
(c) For invertible $\mathbf R$ and positive-definite $\mathsf G'$, $\|\mathbf R^{-1}\Delta\mathbf R\|_{\mathsf G'}=\|\Delta\|_{\mathbf R^{-\top}\mathsf G'\mathbf R^{-1}}$.
\end{lemma}

\begin{proof}
(a) Cyclicity of the trace gives $\langle\mathbf H,\mathbf H'\rangle_{\mathsf G}=\operatorname{tr}((\mathsf G^{1/2}\mathbf H\mathsf G^{-1/2})^\top\mathsf G^{1/2}\mathbf H'\mathsf G^{-1/2})$, the Frobenius inner product of the images under the algebra isomorphism $\mathbf H\mapsto\mathsf G^{1/2}\mathbf H\mathsf G^{-1/2}$, which yields the inner product, submultiplicativity and the norm of $I$. Since $(\mathbf A\mathbf B)^*=\mathbf B^*\mathbf A^*$ and $\mathbf H^{**}=\mathbf H$, $\operatorname{tr}(\mathbf X^*\mathbf H\mathbf Y)=\operatorname{tr}((\mathbf H^*\mathbf X)^*\mathbf Y)$ and $\operatorname{tr}(\mathbf X^*\mathbf Y\mathbf H)=\operatorname{tr}((\mathbf X\mathbf H^*)^*\mathbf Y)$.
(b) For $\mathbf H=\mathbf b\mathbf c^\top$, $\|\mathbf H\|^2_{\mathsf G'}=(\mathbf b^\top\mathsf G'\mathbf b)(\mathbf c^\top\mathsf G'^{-1}\mathbf c)$. If $\mathsf G''=s\mathsf G'$ with $s>0$, the factors $s$ and $s^{-1}$ cancel in the trace form for every matrix. Conversely, fixing $\mathbf c\neq0$ gives $\mathbf b^\top\mathsf G''\mathbf b=s\mathbf b^\top\mathsf G'\mathbf b$ for every $\mathbf b$, with $s=(\mathbf c^\top\mathsf G'^{-1}\mathbf c)/(\mathbf c^\top\mathsf G''^{-1}\mathbf c)>0$, and symmetric matrices with proportional quadratic forms are proportional.
(c) By cyclicity, $\|\mathbf R^{-1}\Delta\mathbf R\|^2_{\mathsf G'}=\operatorname{tr}(\mathbf R\mathsf G'^{-1}\mathbf R^\top\Delta^\top\mathbf R^{-\top}\mathsf G'\mathbf R^{-1}\Delta)$, and $(\mathbf R^{-\top}\mathsf G'\mathbf R^{-1})^{-1}=\mathbf R\mathsf G'^{-1}\mathbf R^\top$.
\end{proof}

\begin{proposition}[Loop defect]\label{prop:llp_loop_defect}
Let $\mathfrak Q_p$ be the loop defect of Eq.~\eqref{eq:pifb_metric_loop_defect}.
(a) $\mathfrak Q_p$ is continuous and nonnegative on $\mathcal T$ and vanishes exactly when $\rho(\mathring\Theta_p)=I$.
(b) A passive frame change leaves $\mathfrak Q_p$ unchanged. Let the reference change by $h\in G$, so that the frames become $U_ih$ and the generators $\operatorname{Ad}_{h^{-1}}\mathring\omega$, and put $\mathbf R=\rho(h)$. The defect computed at the new reference with the matrix $\mathbf R^\top\mathsf G\mathbf R$, which represents the same inner product in the new coordinates, equals $\mathfrak Q_p$ for every generator family. The defect computed with $\mathsf G$ held fixed equals the defect at the old reference computed with $\mathbf R^{-\top}\mathsf G\mathbf R^{-1}$; it equals $\mathfrak Q_p$ for every generator family whenever $\mathbf R$ is a nonzero multiple of a $\mathsf G$-isometry, and, when $\mathrm d\rho(\mathfrak g)=\mathfrak{gl}(\mathbb V)$ and $p$ traverses some reversal pair exactly once, only then.
(c) Put $\mathring\Omega_a=\mathrm d\rho\mathring\omega_{e_a}$ and $\varsigma_p=\sum_a\|\mathring\Omega_a\|_{\mathsf G}$. Then $\sum_a\mathring\Omega_a=\mathrm d\rho(\mathbf C_p\mathring\omega)$ and
\begin{equation}
\Bigl|\mathfrak Q_p-\tfrac12\|\mathrm d\rho(\mathbf C_p\mathring\omega)\|_{\mathsf G}^2\Bigr|
\leq\varsigma_p(e^{\varsigma_p}-1-\varsigma_p)+\tfrac12(e^{\varsigma_p}-1-\varsigma_p)^2
\leq\tfrac{23}{10}\varsigma_p^3\quad\text{if }\varsigma_p\leq1.
\label{eq:llp_defect_expansion}
\end{equation}
If $\mathrm d\rho$ is injective and $\mathfrak g$ carries the inner product pulled back from $\langle\cdot,\cdot\rangle_{\mathsf G}$, then $\varsigma_p\leq r\|\mathring\omega\|$, so $\mathfrak Q_p=\frac12\|\mathbf C_p\mathring\omega\|^2+O(\|\mathring\omega\|^3)$.
\end{proposition}

\begin{proof}
Write $\Delta=I-\rho(\mathring\Theta_p)$.
(a) The exponential, products and $\rho$ are continuous, and the zero set of $\frac12\|\Delta\|_{\mathsf G}^2$ is that of a norm.
(b) A passive change alters neither the generators (Section~\ref{sec:regime_convention}) nor the reference and $\mathsf G$. After the reference change each factor is $\exp(\operatorname{Ad}_{h^{-1}}\mathring\omega_{e_a})=h^{-1}\exp(\mathring\omega_{e_a})h$, so the new reference holonomy is $h^{-1}\mathring\Theta_ph$ and its defect matrix is $\mathbf R^{-1}\Delta\mathbf R$. This is the old matrix written in the coordinates $v'=\mathbf R^{-1}v$, in which the inner product with matrix $\mathsf G$ has the matrix $\mathbf R^\top\mathsf G\mathbf R$. Lemma~\ref{lem:llp_trace_form}(c) with $\mathsf G'=\mathbf R^\top\mathsf G\mathbf R$ returns $\|\Delta\|_{\mathsf G}$, and with $\mathsf G'=\mathsf G$ gives $\|\Delta\|_{\mathbf R^{-\top}\mathsf G\mathbf R^{-1}}$. If $\mathbf R=c\mathbf O$ with $c\neq0$ and $\mathbf O^\top\mathsf G\mathbf O=\mathsf G$, then $\mathbf R^{-\top}\mathsf G\mathbf R^{-1}=c^{-2}\mathsf G$, and Lemma~\ref{lem:llp_trace_form}(b) gives equality. Conversely, suppose that equality holds for every generator family, let $p$ traverse the pair of $e$ exactly once, let $\mathbf H$ be any matrix, choose $\xi\in\mathfrak g$ with $\mathrm d\rho\xi=\mathbf H$, and give the pair of $e$ the generator $t\xi$ and every other pair the generator zero. Then $\rho(\mathring\Theta_p)=\exp(\pm t\mathbf H)$, so $\Delta=\mp t\mathbf H+O(t^2)$; dividing the assumed equality by $t^2$ and letting $t\to0$ gives $\|\mathbf H\|_{\mathbf R^{-\top}\mathsf G\mathbf R^{-1}}=\|\mathbf H\|_{\mathsf G}$ for every $\mathbf H$. Lemma~\ref{lem:llp_trace_form}(b) gives $\mathbf R^{-\top}\mathsf G\mathbf R^{-1}=s\mathsf G$ with $s>0$, so $s^{1/2}\mathbf R$ is a $\mathsf G$-isometry.
(c) The first identity is the linearity of $\mathrm d\rho$ with the reversal convention. Multiplying the exponential series, $\mathbf T_p=\rho(\mathring\Theta_p)-I-\sum_a\mathring\Omega_a$ is the absolutely convergent sum of the ordered products $\mathring\Omega_0^{l_0}\cdots\mathring\Omega_{r-1}^{l_{r-1}}/(l_0!\cdots l_{r-1}!)$ over the multi-indices with $l_0+\cdots+l_{r-1}\geq2$, so Lemma~\ref{lem:llp_trace_form}(a) gives $\|\mathbf T_p\|_{\mathsf G}\leq e^{\varsigma_p}-1-\varsigma_p$. Since $2\mathfrak Q_p=\|\sum_a\mathring\Omega_a+\mathbf T_p\|_{\mathsf G}^2$ and $\|\sum_a\mathring\Omega_a\|_{\mathsf G}\leq\varsigma_p$, expanding the square gives the first inequality. With $e^\varsigma-1-\varsigma\leq\frac12\varsigma^2e^\varsigma$, the right side is at most $\frac12\varsigma_p^3e^{\varsigma_p}+\frac18\varsigma_p^4e^{2\varsigma_p}\leq(\frac e2+\frac{e^2}8)\varsigma_p^3$ for $\varsigma_p\leq1$, and $\frac e2+\frac{e^2}8<2.29$. With the pulled-back inner product, $\|\mathring\Omega_a\|_{\mathsf G}=\|\mathring\omega_{e_a}\|\leq\|\mathring\omega\|$.
\end{proof}

\begin{proposition}[Continuous class functions are blind to unipotent holonomy]\label{prop:llp_class_function}
Let $K\geq2$, let $\mathbb H$ be $\mathrm{GL}(K,\mathbb R)$ or $\mathrm{GL}^+(K)$, and let $f:\mathbb H\to\mathbb R$ be continuous and constant on the conjugacy classes of $\mathbb H$. Then $f(I+\mathbf T)=f(I)$ for every nilpotent $\mathbf T$. In particular no such function vanishes exactly at $I$, since the shear $I+E_{12}\neq I$ is unipotent. Consequently, for $G=\mathrm{GL}(K,\mathbb R)$ or $\mathrm{GL}^+(K)$ in its defining representation and a closed walk that traverses some reversal pair exactly once, a loop penalty that is a continuous function of the holonomy and is unchanged by every passive frame change takes at every unipotent holonomy the value it takes at the identity.
\end{proposition}

\begin{proof}
Since $\mathbf T^K=0$, the kernels $\ker\mathbf T\subseteq\ker\mathbf T^2\subseteq\cdots\subseteq\ker\mathbf T^K=\mathbb R^K$ increase and $\mathbf T$ maps $\ker\mathbf T^{k+1}$ into $\ker\mathbf T^k$. A basis of $\ker\mathbf T$, extended successively to bases of $\ker\mathbf T^2,\ldots,\ker\mathbf T^K$, gives the columns of an invertible $\mathbf B$ with $\mathbf T_0=\mathbf B^{-1}\mathbf T\mathbf B$ strictly upper triangular; replacing the first column by its negative preserves this property, so $\det\mathbf B>0$ may be assumed. For $t>0$ put $\mathbf S_t=\operatorname{diag}(1,t,\ldots,t^{K-1})$. Then $\mathbf B\mathbf S_t\in\mathrm{GL}^+(K)\subseteq\mathbb H$, and $(\mathbf S_t^{-1}\mathbf T_0\mathbf S_t)_{kl}=t^{l-k}(\mathbf T_0)_{kl}\to0$ as $t\to0$, because $(\mathbf T_0)_{kl}=0$ unless $l>k$. Hence $f(I+\mathbf T)=f(I+\mathbf S_t^{-1}\mathbf T_0\mathbf S_t)$ for every $t>0$, and continuity gives $f(I+\mathbf T)=f(I)$. For the last claim, give the traversed pair the link $g'$ or $g'^{-1}$, according to the direction of traversal, and every other link the identity. The holonomy is then $g'$, and by Eq.~\eqref{eq:discrete_holonomy} a passive change at the base vertex replaces it by $gg'g^{-1}$, so the penalty is a continuous class function on $G$.
\end{proof}

\begin{lemma}[Gaussian moments and growth]\label{lem:llp_growth}
(a) For a Gaussian law $\mathcal N(\mu,\Sigma)$ on a Euclidean space of dimension $d$ and every $c\geq0$, $\mathbb Ee^{c\|\xi\|}\leq e^{c\|\mu\|}(2e^{c^2\|\Sigma\|/2})^d$.
(b) The matrix exponential and its first two derivatives at $\mathbf B$, in directions $\mathbf T$ and $\mathbf T'$, have norms at most $e^{\|\mathbf B\|}$, $e^{\|\mathbf B\|}\|\mathbf T\|$ and $e^{\|\mathbf B\|}\|\mathbf T\|\|\mathbf T'\|$.
(c) Every real polynomial $f$ in the entries of the matrices $\exp(\pm\mathrm d\rho\mathring\omega_e)$, $e\in E_+$, is smooth on $\mathcal T$ with $|f|+\|\nabla f\|+\|\nabla^2f\|\leq c_1e^{c_2\|\mathring\omega\|}$. Each loop defect, and hence $\mathcal A_{\mathrm{loop}}$, is such a polynomial.
\end{lemma}

\begin{proof}
(a) Write $\xi=\mu+\Sigma^{1/2}\mathbf b$ with $\mathbf b$ standard Gaussian in an orthonormal basis. Then $\|\xi\|\leq\|\mu\|+\|\Sigma\|^{1/2}\sum_l|b_l|$, the coordinates $b_l$ are independent, and $\mathbb Ee^{c'|b_l|}\leq\mathbb Ee^{c'b_l}+\mathbb Ee^{-c'b_l}=2e^{c'^2/2}$.
(b) The first derivative of $\mathbf B\mapsto\mathbf B^l$ in direction $\mathbf T$ is a sum of $l$ products of $l-1$ factors $\mathbf B$ and one factor $\mathbf T$, and the second derivative is a sum of $l(l-1)$ products of $l-2$ factors $\mathbf B$ and one factor each of $\mathbf T$ and $\mathbf T'$; summing the exponential series gives the bounds.
(c) By (b) and the linearity of $\mathrm d\rho$, every entry of $\exp(\pm\mathrm d\rho\mathring\omega_e)$ and its first two derivatives in $\mathring\omega$ are bounded by $ce^{\|\mathrm d\rho\|\|\mathring\omega\|}$, and the Leibniz rule bounds a monomial of degree $k$ and its first two derivatives by $c'e^{k\|\mathrm d\rho\|\|\mathring\omega\|}$. A loop defect is a quadratic polynomial in the entries of $\rho(\mathring\Theta_p)$, which is a product of the matrices $\exp(\pm\mathrm d\rho\mathring\omega_e)$.
\end{proof}

\begin{proposition}[Normalization and covariance of the generator prior]\label{prop:llp_generator_prior}
Let $G$ be a Lie group, compact or not, and let the data of Eq.~\eqref{eq:pifb_generator_prior} be given.
(a) $\mathcal A_{\mathrm{loop}}$ is finite, continuous and nonnegative, has finite $\nu^{\mathcal T}$-mean, and $0<e^{-\mathbb E_{\nu^{\mathcal T}}\mathcal A_{\mathrm{loop}}}\leq Z_{\mathrm{loop}}\leq1$. If $\rho(G)$ consists of $\mathsf G$-isometries, also $Z_{\mathrm{loop}}\geq\exp(-2n_V\sum_pc_p)$. The prior $P_{\mathrm{gen}}$ is a probability law with a strictly positive continuous density with respect to Lebesgue measure on $\mathcal T$.
(b) Passive frame changes leave $P_{\mathrm{gen}}$ unchanged.
(c) Under the reference change by $h$, with $\mathbf R=\rho(h)$ and $\operatorname{Ad}_{h^{-1}}$ acting on every summand, the image of $P_{\mathrm{gen}}$ is the generator prior at the new reference with anchors $\operatorname{Ad}_{h^{-1}}\mathring\alpha_e$, covariances $\operatorname{Ad}_{h^{-1}}\mathring\Sigma_e\operatorname{Ad}_h$, fiber matrix $\mathbf R^\top\mathsf G\mathbf R$, inner product $\langle\operatorname{Ad}_h\cdot,\operatorname{Ad}_h\cdot\rangle$ on $\mathfrak g$, and the same walks and coefficients; both priors induce the same law of links.
(d) If $\operatorname{Ad}_h$ is an isometry of the inner product on $\mathfrak g$, fixes every anchor and commutes with every covariance, and $\mathbf R$ is a nonzero multiple of a $\mathsf G$-isometry, then $P_{\mathrm{gen}}$ is invariant under $\operatorname{Ad}_{h^{-1}}$.
\end{proposition}

\begin{proof}
(a) Finiteness, continuity and nonnegativity follow from Proposition~\ref{prop:llp_loop_defect}(a). The exponential series and Lemma~\ref{lem:llp_trace_form}(a) give $\|\exp\mathbf H\|_{\mathsf G}\leq n_V^{1/2}+e^{\|\mathbf H\|_{\mathsf G}}-1\leq n_V^{1/2}e^{\|\mathbf H\|_{\mathsf G}}$, so a walk of length $r$ has $\mathfrak Q_p\leq\frac12(n_V^{1/2}+n_V^{r/2}e^{\varsigma_p})^2\leq2n_V^re^{2\varsigma_p}$. Here $\varsigma_p$ is at most the operator norm of $\mathrm d\rho$ times the sum of the norms of the generators on the traversed pairs, each counted with its number of traversals. Under $\nu^{\mathcal T}$ the generators of distinct pairs are independent Gaussian vectors, so Lemma~\ref{lem:llp_growth}(a) gives $\mathbb E_{\nu^{\mathcal T}}e^{2\varsigma_p}<\infty$, and $\mathcal A_{\mathrm{loop}}$ has a finite mean. Since $0<e^{-\mathcal A_{\mathrm{loop}}}\leq1$, $0<Z_{\mathrm{loop}}\leq1$, and Jensen's inequality for the convex function $u\mapsto e^{-u}$ gives $Z_{\mathrm{loop}}\geq e^{-\mathbb E\mathcal A_{\mathrm{loop}}}$. If $\mathbf B=\rho(\mathring\Theta_p)$ satisfies $\mathbf B^\top\mathsf G\mathbf B=\mathsf G$, then $\Delta^\top\mathsf G\Delta=2\mathsf G-\mathbf B^\top\mathsf G-\mathsf G\mathbf B$ and $\operatorname{tr}(\mathsf G^{-1}\mathbf B^\top\mathsf G)=\operatorname{tr}\mathbf B$, so $\mathfrak Q_p=n_V-\operatorname{tr}\mathbf B\leq2n_V$, because $\mathsf G^{1/2}\mathbf B\mathsf G^{-1/2}$ is orthogonal and its eigenvalues have modulus one; hence $\mathcal A_{\mathrm{loop}}\leq2n_V\sum_pc_p$. The density of $P_{\mathrm{gen}}$ is $Z_{\mathrm{loop}}^{-1}e^{-\mathcal A_{\mathrm{loop}}}$ times a product of nondegenerate Gaussian densities.
(b) Passive changes alter neither the generators nor the anchors, covariances, inner product, walks and fiber matrix.
(c) By construction $\operatorname{Ad}_{h^{-1}}$ is a linear isometry from $(\mathfrak g,\langle\cdot,\cdot\rangle)$ onto $(\mathfrak g,\langle\operatorname{Ad}_h\cdot,\operatorname{Ad}_h\cdot\rangle)$. It carries the normalized Lebesgue measure to the normalized Lebesgue measure and $\mathcal N(\mathring\alpha_e,\mathring\Sigma_e)$ to the Gaussian law with mean $\operatorname{Ad}_{h^{-1}}\mathring\alpha_e$ and covariance $\operatorname{Ad}_{h^{-1}}\mathring\Sigma_e\operatorname{Ad}_h$, since for every $\beta$ the new pairing of $\beta$ with $\operatorname{Ad}_{h^{-1}}\mathring\omega_e$ is $\langle\operatorname{Ad}_h\beta,\mathring\omega_e\rangle$, of variance $\langle\operatorname{Ad}_h\beta,\mathring\Sigma_e\operatorname{Ad}_h\beta\rangle$. When $\mathfrak g$ carries the pulled-back inner product, $\mathrm d\rho(\operatorname{Ad}_h\xi)=\mathbf R\mathrm d\rho(\xi)\mathbf R^{-1}$ and Lemma~\ref{lem:llp_trace_form}(c) identify the new inner product with the pullback for $\mathbf R^\top\mathsf G\mathbf R$. By Proposition~\ref{prop:llp_loop_defect}(b) the new defect of $\operatorname{Ad}_{h^{-1}}\mathring\omega$ equals $\mathfrak Q_p(\mathring\omega)$, so the transformed action composed with $\operatorname{Ad}_{h^{-1}}$ is $\mathcal A_{\mathrm{loop}}$ and the normalizer is unchanged. The links agree because $U_ih\exp(\operatorname{Ad}_{h^{-1}}\mathring\omega_e)h^{-1}U_j^{-1}=U_i\exp(\mathring\omega_e)U_j^{-1}$.
(d) Under these hypotheses the data in (c) are the original anchors, covariances, inner product, walks and coefficients, with the fiber matrix multiplied by a positive scalar, which changes neither the trace form nor the defect by Lemma~\ref{lem:llp_trace_form}(b).
\end{proof}

\subsection{Latent generators in the evidence bound}\label{app:llp_decomposition}

Assume the setting of Eq.~\eqref{eq:pifb_latent_population_joint}: the map $\mathring\omega\mapsto\Theta(\mathring\omega)$ of Eq.~\eqref{eq:pifb_latent_link_chart} is continuous, and the held-link prior and record kernels are measurable in the link values, so that $\mathring\omega\mapsto\mathbb P_{\Theta(\mathring\omega)}$ is a probability kernel. Recognition laws are disintegrated as $Q(d\mathring\omega,dy)=R(d\mathring\omega)Q_{\mathring\omega}(dy)$.

\begin{proposition}[Normalization, evidence and conditional decomposition]\label{prop:llp_decomposition}
(a) For every finite record family, cyclic or not, the joint law~\eqref{eq:pifb_latent_population_joint} is a probability law.
(b) The conditional evidence $\mathring\omega\mapsto p(o\mid\mathring\omega)$ is measurable, and the evidence of the joint is $p(o)=\int p(o\mid\mathring\omega)P_{\mathrm{gen}}(d\mathring\omega)$. It is positive exactly when $p(o\mid\cdot)$ is positive on a set of positive Lebesgue measure. If $\ell_o\leq c_{\mathfrak A}$, then $p(o\mid\mathring\omega)\leq c_{\mathfrak A}$ for every $\mathring\omega$ and $p(o)\leq c_{\mathfrak A}$.
(c) Suppose that $\ell_o\leq c_{\mathfrak A}$ and $p(o\mid\mathring\omega)>0$ for every $\mathring\omega$. Then the posterior disintegrates as $\Pi^{\mathcal T}(d\mathring\omega)\Pi_{\mathring\omega}(dy)$ with $\Pi^{\mathcal T}=p(o)^{-1}p(o\mid\cdot)P_{\mathrm{gen}}$, and every law $Q=R(d\mathring\omega)Q_{\mathring\omega}(dy)$ satisfies
\begin{equation}
\mathcal F_V[Q;X,o]=D_{\mathrm{KL}}(R\|P_{\mathrm{gen}})+\int\mathcal F_V[Q_{\mathring\omega};X,\Theta(\mathring\omega),o]R(d\mathring\omega)
\quad\text{in }(-\infty,+\infty],
\label{eq:llp_decomposition}
\end{equation}
with a measurable integrand bounded below by $-\log c_{\mathfrak A}$. If $D_{\mathrm{KL}}(R\|\nu^{\mathcal T})<\infty$ and $\mathbb E_R\mathcal A_{\mathrm{loop}}<\infty$, then $D_{\mathrm{KL}}(R\|P_{\mathrm{gen}})=D_{\mathrm{KL}}(R\|\nu^{\mathcal T})+\mathbb E_R\mathcal A_{\mathrm{loop}}+\log Z_{\mathrm{loop}}$.
\end{proposition}

\begin{proof}
(a) $P_{\mathrm{gen}}$ is a probability law by Proposition~\ref{prop:llp_generator_prior}(a), and for every link family the held-link law has unit mass, by integrating the records and then the prior as in Eq.~\eqref{eq:population_joint}. Tonelli's theorem, integrating records, then matter, then generators, gives unit mass; cycles in the record hypergraph do not enter, because every record kernel is normalized conditional on the whole latent state.
(b) Joint measurability of $\ell_o$ and Tonelli's theorem give the measurability of $p(o\mid\cdot)$ and the evidence formula. Since $P_{\mathrm{gen}}$ has a strictly positive Lebesgue density, the integral of the nonnegative function $p(o\mid\cdot)$ is positive exactly when that function is positive on a set of positive Lebesgue measure. The bounds follow by integrating $\ell_o\leq c_{\mathfrak A}$ against probability laws.
(c) The measure $\ell_o(y;\Theta(\mathring\omega))P^Y_{\mathring\omega}(dy)P_{\mathrm{gen}}(d\mathring\omega)$ equals $p(o\mid\mathring\omega)\Pi_{\mathring\omega}(dy)P_{\mathrm{gen}}(d\mathring\omega)$, and $\mathring\omega\mapsto\Pi_{\mathring\omega}$ is a kernel by Tonelli's theorem; dividing by $p(o)\in(0,c_{\mathfrak A}]$ gives the disintegration. The chain rule~\eqref{eq:llp_chain_rule} gives $D_{\mathrm{KL}}(Q\|\Pi)=D_{\mathrm{KL}}(R\|\Pi^{\mathcal T})+\int D_{\mathrm{KL}}(Q_{\mathring\omega}\|\Pi_{\mathring\omega})R(d\mathring\omega)$ with a measurable integrand. If $D_{\mathrm{KL}}(R\|P_{\mathrm{gen}})<\infty$, then $R\ll P_{\mathrm{gen}}$, and since $d\Pi^{\mathcal T}/dP_{\mathrm{gen}}=p(o\mid\cdot)/p(o)$ is positive and finite, $R\ll\Pi^{\mathcal T}$ with $\log(dR/d\Pi^{\mathcal T})=\log(dR/dP_{\mathrm{gen}})-\log p(o\mid\mathring\omega)+\log p(o)$. The first term has an $R$-integrable negative part and the second is bounded below by $-\log c_{\mathfrak A}$, so the integral splits; adding the conditional term and $-\log p(o)$ gives Eq.~\eqref{eq:llp_decomposition}, because $\mathcal F_V[Q_{\mathring\omega};X,\Theta(\mathring\omega),o]=-\log p(o\mid\mathring\omega)+D_{\mathrm{KL}}(Q_{\mathring\omega}\|\Pi_{\mathring\omega})\geq-\log c_{\mathfrak A}$. If $D_{\mathrm{KL}}(R\|P_{\mathrm{gen}})=+\infty$, then either $R\not\ll P_{\mathrm{gen}}$, hence $R\not\ll\Pi^{\mathcal T}$ because the two laws have the same null sets, or $R\ll P_{\mathrm{gen}}$ and $\Pi^{\mathcal T}\leq(c_{\mathfrak A}/p(o))P_{\mathrm{gen}}$ gives $dR/d\Pi^{\mathcal T}\geq(p(o)/c_{\mathfrak A})dR/dP_{\mathrm{gen}}$. In both cases $D_{\mathrm{KL}}(R\|\Pi^{\mathcal T})=+\infty$, and both sides of Eq.~\eqref{eq:llp_decomposition} are $+\infty$. The last claim integrates $\log(dR/dP_{\mathrm{gen}})=\log(dR/d\nu^{\mathcal T})+\mathcal A_{\mathrm{loop}}+\log Z_{\mathrm{loop}}$ against $R$.
\end{proof}

\begin{proposition}[Gaussian recognition of the generators]\label{prop:llp_gaussian_recognition}
Let $R=\mathcal N(\mathring\mu,\mathring\Lambda^{-1})$ with $\mathring\Lambda\succ0$, and let $\mathring\Sigma$ be the block-diagonal prior covariance. Then
\begin{equation}
D_{\mathrm{KL}}(R\|\nu^{\mathcal T})=\tfrac12\bigl[\operatorname{tr}(\mathring\Sigma^{-1}\mathring\Lambda^{-1})+\langle\mathring\mu-\mathring\alpha,\mathring\Sigma^{-1}(\mathring\mu-\mathring\alpha)\rangle-\dim\mathcal T+\log\det\mathring\Sigma+\log\det\mathring\Lambda\bigr],
\label{eq:llp_gaussian_kl}
\end{equation}
$\mathbb E_Re^{c\|\mathring\omega\|}<\infty$ for every $c\geq0$, and $\mathbb E_R\mathcal A_{\mathrm{loop}}<\infty$. Hence, under the hypotheses of Proposition~\ref{prop:llp_decomposition}(c), the link sector of Eq.~\eqref{eq:pifb_latent_collective_vfe} is finite and equals $D_{\mathrm{KL}}(R\|P_{\mathrm{gen}})$.
\end{proposition}

\begin{proof}
For a self-adjoint $\mathbf B$, $\mathbb E_R\langle\mathring\omega-\mathring\alpha,\mathbf B(\mathring\omega-\mathring\alpha)\rangle=\operatorname{tr}(\mathbf B\mathring\Lambda^{-1})+\langle\mathring\mu-\mathring\alpha,\mathbf B(\mathring\mu-\mathring\alpha)\rangle$, and the $R$-expectation of the difference of the two Gaussian log densities gives Eq.~\eqref{eq:llp_gaussian_kl}. The moment bound is Lemma~\ref{lem:llp_growth}(a), and Lemma~\ref{lem:llp_growth}(c) gives $\mathbb E_R\mathcal A_{\mathrm{loop}}\leq c_1\mathbb E_Re^{c_2\|\mathring\omega\|}$. The last claim is Proposition~\ref{prop:llp_decomposition}(c).
\end{proof}

\subsection{Held links as a limit, point masses and conditional marginals}\label{app:llp_held}

Hold the anchor $\mathring\alpha$, the record and the structural data, and let the prior covariance be $\mathring\Sigma=\mathring\sigma^2\mathring\Sigma_0$ with $\mathring\sigma>0$ and a block-diagonal $\mathring\Sigma_0\succ0$. Write $\nu_{\mathring\sigma}$, $P_{\mathring\sigma}$, $Z_{\mathring\sigma}$ and $\mathcal F_{\mathring\sigma}$ for the base law, generator prior, normalizer and latent-link free energy at width $\mathring\sigma$. For a class $\mathcal Q'$ of laws on the matter space put $\mathbf F_{Q'}(\mathring\omega)=\mathcal F_V[Q';X,\Theta(\mathring\omega),o]$ and $\mathbf F(\mathring\omega)=\inf_{Q'\in\mathcal Q'}\mathbf F_{Q'}(\mathring\omega)$.

\begin{theorem}[Held links are the zero-width limit]\label{thm:llp_held_limit}
Assume (i) $\ell_o\leq c_{\mathfrak A}$ and $p(o\mid\mathring\omega)>0$ for every $\mathring\omega$; (ii) $\mathbf F$ is finite and lower semicontinuous at $\mathring\alpha$; and (iii) a subclass $\mathcal Q'_0\subseteq\mathcal Q'$ with $\inf_{Q'\in\mathcal Q'_0}\mathbf F_{Q'}(\mathring\alpha)=\mathbf F(\mathring\alpha)$ consists of laws for which $\mathbf F_{Q'}$ is continuous at $\mathring\alpha$ and bounded above by $c_{Q'}e^{c_{Q'}\|\mathring\omega\|}$. Then, as $\mathring\sigma\to0$, the infimum of $\mathcal F_{\mathring\sigma}$ over the product laws $R\otimes Q'$ with $Q'\in\mathcal Q'$ and $R$ in any class of laws containing $\nu_{\mathring\sigma}$, and its infimum over all laws $R(d\mathring\omega)Q_{\mathring\omega}(dy)$ whose kernel takes values in $\mathcal Q'$, both converge to $\mathbf F(\mathring\alpha)$, the infimum over $\mathcal Q'$ of the held-link free energy at the anchor links $U_i\exp(\mathring\alpha_e)U_j^{-1}$.
\end{theorem}

\begin{proof}
Write $\mathring\omega=\mathring\alpha+\mathring\sigma\mathring\Sigma_0^{1/2}\mathbf b$ under $\nu_{\mathring\sigma}$, with $\mathbf b$ standard Gaussian. Since $e^{-\mathcal A_{\mathrm{loop}}}\in(0,1]$ is continuous, bounded convergence gives $Z_{\mathring\sigma}\to e^{-\mathcal A_{\mathrm{loop}}(\mathring\alpha)}>0$, so $P_{\mathring\sigma}\{\|\mathring\omega-\mathring\alpha\|\geq\epsilon\}\leq\nu_{\mathring\sigma}\{\|\mathring\omega-\mathring\alpha\|\geq\epsilon\}/Z_{\mathring\sigma}\to0$ for every $\epsilon>0$. By Proposition~\ref{prop:llp_decomposition}(c), every law $Q=R(d\mathring\omega)Q_{\mathring\omega}(dy)$ has $\mathcal F_{\mathring\sigma}[Q]=D_{\mathrm{KL}}(R\|P_{\mathring\sigma})+\int\mathbf F_{Q_{\mathring\omega}}(\mathring\omega)R(d\mathring\omega)$ with $\mathbf F_{Q_{\mathring\omega}}\geq-\log c_{\mathfrak A}$.

For the upper bound, let $Q'\in\mathcal Q'_0$. The law $\nu_{\mathring\sigma}\otimes Q'$ is admissible, and $D_{\mathrm{KL}}(\nu_{\mathring\sigma}\|P_{\mathring\sigma})=\mathbb E_{\nu_{\mathring\sigma}}\mathcal A_{\mathrm{loop}}+\log Z_{\mathring\sigma}$. For $\mathring\sigma\leq1$, Lemma~\ref{lem:llp_growth}(c) dominates $\mathcal A_{\mathrm{loop}}(\mathring\alpha+\mathring\sigma\mathring\Sigma_0^{1/2}\mathbf b)$ by $c_1e^{c_2(\|\mathring\alpha\|+\|\mathring\Sigma_0\|^{1/2}\|\mathbf b\|)}$, which is integrable by Lemma~\ref{lem:llp_growth}(a), so $\mathbb E_{\nu_{\mathring\sigma}}\mathcal A_{\mathrm{loop}}\to\mathcal A_{\mathrm{loop}}(\mathring\alpha)$ and the divergence tends to zero. Hypothesis (iii), the lower bound $-\log c_{\mathfrak A}$ and the same domination give $\mathbb E_{\nu_{\mathring\sigma}}\mathbf F_{Q'}\to\mathbf F_{Q'}(\mathring\alpha)$. Hence the $\limsup$ of the product infimum is at most $\mathbf F_{Q'}(\mathring\alpha)$, and the infimum over $\mathcal Q'_0$ gives at most $\mathbf F(\mathring\alpha)$.

For the lower bound, let $c_0>0$. By (ii) there is $\epsilon>0$ with $\mathbf F>\mathbf F(\mathring\alpha)-c_0$ on the ball $B_\epsilon=\{\|\mathring\omega-\mathring\alpha\|<\epsilon\}$. The bounded function $\mathbf q=(\mathbf F(\mathring\alpha)-c_0)\mathbf 1_{B_\epsilon}-(\log c_{\mathfrak A})\mathbf 1_{B_\epsilon^c}$ satisfies $\mathbf q(\mathring\omega)\leq\mathbf F_{Q_{\mathring\omega}}(\mathring\omega)$ for every kernel with values in $\mathcal Q'$. For every $R$ with $D_{\mathrm{KL}}(R\|P_{\mathring\sigma})<\infty$, the law $P^{\mathbf q}\propto e^{-\mathbf q}P_{\mathring\sigma}$ gives $0\leq D_{\mathrm{KL}}(R\|P^{\mathbf q})=D_{\mathrm{KL}}(R\|P_{\mathring\sigma})+\mathbb E_R\mathbf q+\log\mathbb E_{P_{\mathring\sigma}}e^{-\mathbf q}$, hence $\mathcal F_{\mathring\sigma}[Q]\geq D_{\mathrm{KL}}(R\|P_{\mathring\sigma})+\mathbb E_R\mathbf q\geq-\log\mathbb E_{P_{\mathring\sigma}}e^{-\mathbf q}$; the bound is trivial when the divergence is infinite. Since $P_{\mathring\sigma}(B_\epsilon^c)\to0$, the right side tends to $\mathbf F(\mathring\alpha)-c_0$. The kernel infimum is at most the product infimum, so both converge to $\mathbf F(\mathring\alpha)$.
\end{proof}

\begin{proposition}[Point-mass recognition]\label{prop:llp_point_mass}
For every $\mathring\mu$, $D_{\mathrm{KL}}(\delta_{\mathring\mu}\|P_{\mathrm{gen}})=+\infty$. For $R=\mathcal N(\mathring\mu,\mathring\Lambda^{-1})$,
$D_{\mathrm{KL}}(R\|P_{\mathrm{gen}})\geq\frac12(\log\det\mathring\Lambda-\log\det\mathring\Sigma^{-1}-\dim\mathcal T)+\log Z_{\mathrm{loop}}$, which tends to $+\infty$ as $\mathring\Lambda^{-1}\to0$ with the prior held. Under the hypotheses of Proposition~\ref{prop:llp_decomposition}(c), the free energy of $\delta_{\mathring\mu}\otimes Q_Y$ is $+\infty$ and that of $R\otimes Q_Y$ tends to $+\infty$.
\end{proposition}

\begin{proof}
$P_{\mathrm{gen}}$ has a Lebesgue density, so $\{\mathring\mu\}$ is $P_{\mathrm{gen}}$-null and $\delta_{\mathring\mu}\not\ll P_{\mathrm{gen}}$. For Gaussian $R$, Propositions~\ref{prop:llp_decomposition}(c) and~\ref{prop:llp_gaussian_recognition} with $\mathcal A_{\mathrm{loop}}\geq0$ give $D_{\mathrm{KL}}(R\|P_{\mathrm{gen}})\geq D_{\mathrm{KL}}(R\|\nu^{\mathcal T})+\log Z_{\mathrm{loop}}$; in Eq.~\eqref{eq:llp_gaussian_kl} the trace and quadratic terms are nonnegative, while $\log\det\mathring\Lambda\geq\dim\mathcal T\log\lambda_{\min}(\mathring\Lambda)\to+\infty$. The free-energy statements follow from Eq.~\eqref{eq:llp_decomposition}, whose matter integrand is bounded below by $-\log c_{\mathfrak A}$.
\end{proof}

\begin{proposition}[Gaussian conditionals with a non-Gaussian marginal]\label{prop:llp_nongaussian_marginal}
Let $G=\mathrm{GL}^+(1)$ act on $\mathbb R$ by multiplication, with trivial frames, $\mathcal A_{\mathrm{loop}}=0$ and $\nu^{\mathcal T}=\mathcal N(0,1)$ on $\mathfrak g=\mathbb R$. Let two scalar agents have $k_j\sim\mathcal N(0,1)$ and $k_i\mid k_j,\mathring\omega\sim\mathcal N(e^{\mathring\omega}k_j,1)$ for the copy $e:j\to i$, and let there be no record. Every held-link law of $(k_i,k_j)$ is Gaussian, while the $k_i$-marginal of the latent-link law satisfies $\mathbb Ek_i^4=3(1+2e^2+e^8)>3(1+e^2)^2=3(\mathbb Ek_i^2)^2$ and is therefore not Gaussian.
\end{proposition}

\begin{proof}
Given $\mathring\omega$, $(k_i,k_j)$ is centered Gaussian with $\mathbb E[k_i^2\mid\mathring\omega]=1+e^{2\mathring\omega}$, and a centered Gaussian variable has $\mathbb Ek^4=3(\mathbb Ek^2)^2$. With $\mathbb Ee^{t\mathring\omega}=e^{t^2/2}$, $\mathbb Ek_i^2=1+e^2$ and $\mathbb Ek_i^4=3\mathbb E(1+e^{2\mathring\omega})^2=3(1+2e^2+e^8)$. The inequality is strict Jensen, since $e^{2\mathring\omega}$ is not almost surely constant. The marginal is centered, and a centered Gaussian law would give equality.
\end{proof}

\subsection{The link field equation}\label{app:llp_field}

For square matrices write $\langle\mathbf S,\mathbf T\rangle=\operatorname{tr}(\mathbf S^\top\mathbf T)$, and let $(\mathrm d\rho)^\top$ be the adjoint of $\mathrm d\rho$ for the inner product on $\mathfrak g$ and this pairing.

\begin{lemma}[Derivative of the exponential chart]\label{lem:llp_exp_derivative}
Let $\mathbf L(\mathbf B,\mathbf T)=\int_0^1e^{u\mathbf B}\mathbf Te^{(1-u)\mathbf B}du$.
(a) $\|e^{\mathbf B+\mathbf T}-e^{\mathbf B}-\mathbf L(\mathbf B,\mathbf T)\|\leq\frac12\|\mathbf T\|^2e^{\|\mathbf B\|+\|\mathbf T\|}$ and $\|\mathbf L(\mathbf B,\mathbf T)\|\leq e^{\|\mathbf B\|}\|\mathbf T\|$, so $\mathbf T\mapsto\mathbf L(\mathbf B,\mathbf T)$ is the Fr\'echet derivative of the exponential at $\mathbf B$~\cite{Higham2008}.
(b) $\langle\mathbf S,\mathbf L(\mathbf B,\mathbf T)\rangle=\langle\mathbf L(\mathbf B^\top,\mathbf S),\mathbf T\rangle$ and $\langle\mathbf S,\mathbf L(\mathbf B,\mathbf T)\rangle_{\mathsf G}=\langle\mathbf L(\mathbf B^*,\mathbf S),\mathbf T\rangle_{\mathsf G}$.
(c) If $f$ is differentiable at $e^{\mathbf B}$, then $\mathbf B\mapsto f(e^{\mathbf B})$ has gradient $\mathbf L(\mathbf B^\top,\nabla f(e^{\mathbf B}))$ at $\mathbf B$, and $\xi\mapsto f(\exp(\mathrm d\rho\xi))$ has gradient $(\mathrm d\rho)^\top\mathbf L((\mathrm d\rho\xi)^\top,\nabla f(\exp(\mathrm d\rho\xi)))$ at $\xi$.
\end{lemma}

\begin{proof}
For square $\mathbf B,\mathbf B'$, $e^{\mathbf B'}-e^{\mathbf B}=\int_0^1e^{u\mathbf B'}(\mathbf B'-\mathbf B)e^{(1-u)\mathbf B}du$, because the integrand is the $u$-derivative of $e^{u\mathbf B'}e^{(1-u)\mathbf B}$; with $\|e^{u\mathbf B}\|\leq e^{u\|\mathbf B\|}$ this gives $\|e^{\mathbf B'}-e^{\mathbf B}\|\leq\|\mathbf B'-\mathbf B\|e^{\max(\|\mathbf B'\|,\|\mathbf B\|)}$. With $\mathbf B'=\mathbf B+\mathbf T$ the remainder in (a) is $\int_0^1(e^{u(\mathbf B+\mathbf T)}-e^{u\mathbf B})\mathbf Te^{(1-u)\mathbf B}du$, whose integrand is bounded by $u\|\mathbf T\|^2e^{\|\mathbf B\|+\|\mathbf T\|}$; the bound on $\mathbf L$ is immediate. For (b), cyclicity and $(e^{u\mathbf B})^\top=e^{u\mathbf B^\top}$ give $\operatorname{tr}(\mathbf S^\top e^{u\mathbf B}\mathbf Te^{(1-u)\mathbf B})=\operatorname{tr}((e^{u\mathbf B^\top}\mathbf Se^{(1-u)\mathbf B^\top})^\top\mathbf T)$, and integration over $u$ gives the first identity; the second follows in the same way from Lemma~\ref{lem:llp_trace_form}(a) and $(e^{u\mathbf B})^*=e^{u\mathbf B^*}$. Part (c) is the chain rule with (a) and (b), the derivative of $\xi\mapsto\exp(\mathrm d\rho\xi)$ in direction $\eta$ being $\mathbf L(\mathrm d\rho\xi,\mathrm d\rho\eta)$.
\end{proof}

\begin{lemma}[Gaussian derivative identities]\label{lem:llp_gaussian_gradients}
Let $f$ be twice continuously differentiable on a Euclidean space with $|f|+\|\nabla f\|+\|\nabla^2f\|\leq c_1e^{c_2\|\xi\|}$. Then $(\mu,S)\mapsto\mathbb E_{\mathcal N(\mu,S)}f$ is continuously differentiable on the product of the space and the open cone of positive-definite operators, and for self-adjoint $\Delta$
\begin{equation}
\nabla_\mu\mathbb E_{\mathcal N(\mu,S)}f=\mathbb E_{\mathcal N(\mu,S)}\nabla f,\qquad
\frac{d}{dt}\mathbb E_{\mathcal N(\mu,S+t\Delta)}f\Big|_{t=0}=\tfrac12\operatorname{tr}\bigl(\Delta\mathbb E_{\mathcal N(\mu,S)}\nabla^2f\bigr).
\label{eq:llp_bonnet_price}
\end{equation}
\end{lemma}

These are the identities used for variational Gaussian approximations~\cite{OpperArchambeau2009}.

\begin{proof}
All expectations below are finite by Lemma~\ref{lem:llp_growth}(a). For a continuously differentiable $\varphi$ with $|\varphi|+\|\nabla\varphi\|\leq c_1e^{c_2\|\xi\|}$, Gaussian integration by parts gives
\begin{equation}
\mathbb E_{\mathcal N(\mu,S)}[(\xi-\mu)\varphi(\xi)]=S\mathbb E_{\mathcal N(\mu,S)}\nabla\varphi.
\label{eq:llp_stein}
\end{equation}
Indeed, for a standard Gaussian vector $\mathbf b$ in an orthonormal basis, integration by parts in the $l$-th coordinate against the standard normal density, whose derivative is $-b_l$ times the density, has vanishing boundary terms and gives $\mathbb E[b_l\varphi'(\mathbf b)]=\mathbb E[\partial_l\varphi'(\mathbf b)]$, and $\varphi'(\mathbf b)=\varphi(\mu+S^{1/2}\mathbf b)$ has gradient $S^{1/2}\nabla\varphi$. The mean derivative follows by differentiating $\mathbb Ef(\mu+S^{1/2}\mathbf b)$ under the expectation, with domination by $c_1e^{c_2(\|\mu\|+\|\eta\|+\|S\|^{1/2}\|\mathbf b\|)}$ along $\mu+t\eta$, $|t|\leq1$. For the covariance, $S_t=S+t\Delta$ stays in a compact subset of the cone for small $|t|$, so the $t$-derivative of the Gaussian density, $\frac12[\langle\xi-\mu,S_t^{-1}\Delta S_t^{-1}(\xi-\mu)\rangle-\operatorname{tr}(S_t^{-1}\Delta)]$ times the density, is dominated uniformly by a Gaussian density times a quadratic polynomial, and the derivative passes under the integral:
\begin{equation*}
\frac{d}{dt}\mathbb Ef\Big|_{t=0}=\tfrac12\mathbb E[f\langle\xi-\mu,S^{-1}\Delta S^{-1}(\xi-\mu)\rangle]-\tfrac12\operatorname{tr}(S^{-1}\Delta)\mathbb Ef.
\end{equation*}
Applying Eq.~\eqref{eq:llp_stein} to $\varphi=f\langle\eta,\xi-\mu\rangle$ for fixed $\eta$, and then to the partial derivatives of $f$, gives $\mathbb E[f(\xi-\mu)\otimes(\xi-\mu)]=S(\mathbb E\nabla^2f)S+S\mathbb Ef$, so the first term equals $\frac12\operatorname{tr}(\Delta\mathbb E\nabla^2f)+\frac12\operatorname{tr}(S^{-1}\Delta)\mathbb Ef$, which is Eq.~\eqref{eq:llp_bonnet_price}. Both derivatives are continuous in $(\mu,S)$ by dominated convergence in the representation $\xi=\mu+S^{1/2}\mathbf b$, since $S\mapsto S^{1/2}$ is continuous on the cone.
\end{proof}

Let $\mathring\Xi$ be the space of self-adjoint operators on $\mathcal T$ that preserve each summand of a fixed orthogonal decomposition of $\mathcal T$: a single summand gives the full family, the edge summands the edge-block family, and the lines of an orthonormal basis the diagonal family. Let $\operatorname{pr}_{\mathring\Xi}$ be the orthogonal projection onto $\mathring\Xi$ in the trace inner product; it keeps the diagonal blocks of the decomposition. The inverse of a positive-definite element of $\mathring\Xi$ lies in $\mathring\Xi$.

\begin{proposition}[Link field equation]\label{prop:llp_field_equation}
Assume $\ell_o\leq c_{\mathfrak A}$ and $p(o\mid\mathring\omega)>0$ for every $\mathring\omega$, a $\sigma$-finite matter reference $\nu^Y$ with respect to which every $P^Y_{\mathring\omega}$ has a density, and a matter recognition law $Q_Y$ with density $q_Y$ and $\mathbb E_{Q_Y}|\log q_Y|<\infty$. Let $H_o(\mathring\omega,y)=-\log p(o,y\mid X,\Theta(\mathring\omega))$, the negative logarithm of the held-link joint density with respect to $\nu^Y\otimes\nu^O$, be twice continuously differentiable in $\mathring\omega$ for every $y$, with
\begin{equation}
\mathbb E_{Q_Y}\sup_{\|\mathring\omega\|\leq r}\bigl(|H_o|+\|\nabla H_o\|+\|\nabla^2H_o\|\bigr)(\mathring\omega,Y)\leq c_1e^{c_2r}\qquad\text{for every }r\geq0.
\label{eq:llp_matter_domination}
\end{equation}
Restrict recognition to $R\otimes Q_Y$ with $R=\mathcal N(\mathring\mu,\mathring\Lambda^{-1})$ and $\mathring\Lambda$ positive definite in $\mathring\Xi$, and put $\mathbf U=\mathcal A_{\mathrm{loop}}+\mathbb E_{Q_Y}H_o(\cdot,Y)$ and $\mathbf V(\mathring\omega)=\frac12\langle\mathring\omega-\mathring\alpha,\mathring\Sigma^{-1}(\mathring\omega-\mathring\alpha)\rangle+\mathbf U(\mathring\omega)$.
(a) $\mathbf U$ is twice continuously differentiable with $|\mathbf U|+\|\nabla\mathbf U\|+\|\nabla^2\mathbf U\|\leq c_3e^{c_4\|\mathring\omega\|}$, and
\begin{equation}
\mathcal F_V[R\otimes Q_Y;X,o]=D_{\mathrm{KL}}(R\|\nu^{\mathcal T})+\mathbb E_R\mathbf U+\log Z_{\mathrm{loop}}+\mathbb E_{Q_Y}\log q_Y
\label{eq:llp_field_free_energy}
\end{equation}
is finite and continuously differentiable in $(\mathring\mu,\mathring\Lambda)$.
(b) $(\mathring\mu,\mathring\Lambda)$ is stationary exactly when $\mathring\Sigma^{-1}(\mathring\mu-\mathring\alpha)+\mathbb E_R\nabla\mathcal A_{\mathrm{loop}}=-\mathbb E_{R\otimes Q_Y}\nabla H_o$ and $\mathring\Lambda=\operatorname{pr}_{\mathring\Xi}(\mathring\Sigma^{-1}+\mathbb E_R\nabla^2\mathcal A_{\mathrm{loop}}+\mathbb E_{R\otimes Q_Y}\nabla^2H_o)$. For the full family these are Eq.~\eqref{eq:latent_link_stationarity}.
(c) If $H_o$ reads $\mathring\omega_e$ only through $M_e=\exp(\mathrm d\rho\mathring\omega_e)$, as a differentiable function on an open set of matrices containing the represented twists, and $\partial H_o/\partial M_e$ is $Q_Y$-integrable at every $\mathring\omega$, then with $\mathbf J_e=-\mathbb E_{Q_Y}\partial H_o/\partial M_e$ the $e$ component of the mean equation is Eq.~\eqref{eq:latent_link_current}.
(d) The mean equation is $\mathbb E_R\nabla\mathbf V=0$, and $\mathbb E_{\mathcal N(\mathring\mu,\mathring\Lambda^{-1})}\nabla\mathbf V\to\nabla\mathbf V(\mathring\mu)$ as $\mathring\Lambda^{-1}\to0$, uniformly for $\mathring\mu$ in compact sets; every limit point of solutions as $\mathring\Lambda^{-1}\to0$ satisfies $\nabla\mathbf V=0$. If $\mathbf U$ is bounded below, $R^\star\propto e^{-\mathbf U}\nu^{\mathcal T}$ is the unique minimizer of $\mathcal F_V[R\otimes Q_Y;X,o]$ over all laws $R$ on $\mathcal T$, and its log density is $-\mathbf V$ up to a constant.
\end{proposition}

\begin{proof}
(a) By Lemma~\ref{lem:llp_growth}(c), $\mathcal A_{\mathrm{loop}}$ has the stated growth. On each ball $\|\mathring\omega\|\leq r$, Eq.~\eqref{eq:llp_matter_domination} dominates $H_o$ and its first two derivatives by a $Q_Y$-integrable function, so differentiation passes under $\mathbb E_{Q_Y}$, and $\mathbb E_{Q_Y}H_o$ is twice continuously differentiable with derivatives bounded by $c_1e^{c_2r}$ there. At each $\mathring\omega$ the held-link posterior has the density $e^{-H_o(\mathring\omega,\cdot)}/p(o\mid\mathring\omega)$ with respect to $\nu^Y$ and $\mathbb E_{Q_Y}|H_o(\mathring\omega,Y)|<\infty$, so Proposition~\ref{prop:llp_evidence}(c) gives $\mathcal F_V[Q_Y;X,\Theta(\mathring\omega),o]=\mathbb E_{Q_Y}[\log q_Y+H_o(\mathring\omega,Y)]$. Proposition~\ref{prop:llp_decomposition}(c) with $Q_{\mathring\omega}=Q_Y$, Proposition~\ref{prop:llp_gaussian_recognition} and Fubini's theorem, justified by $\mathbb E_R\mathbb E_{Q_Y}|H_o|\leq c_1\mathbb E_Re^{c_2\|\mathring\omega\|}<\infty$, give Eq.~\eqref{eq:llp_field_free_energy}, which is finite by Lemma~\ref{lem:llp_growth}(a). The closed form~\eqref{eq:llp_gaussian_kl} is smooth, $\mathbb E_R\mathbf U$ is continuously differentiable in the mean and the covariance $\mathring\Lambda^{-1}$ by Lemma~\ref{lem:llp_gaussian_gradients}, and inversion is a diffeomorphism of the positive-definite cone of $\mathring\Xi$.
(b) By Eq.~\eqref{eq:llp_gaussian_kl} and Lemma~\ref{lem:llp_gaussian_gradients}, the derivative of Eq.~\eqref{eq:llp_field_free_energy} in the mean along $\eta$ is $\langle\mathring\Sigma^{-1}(\mathring\mu-\mathring\alpha)+\mathbb E_R\nabla\mathbf U,\eta\rangle$, and, because $d\log\det S=\operatorname{tr}(S^{-1}\Delta)$, its derivative in the covariance $S=\mathring\Lambda^{-1}$ along $\Delta\in\mathring\Xi$ is $\frac12\operatorname{tr}(\Delta(\mathring\Sigma^{-1}-\mathring\Lambda+\mathbb E_R\nabla^2\mathbf U))$. Stationarity in $(\mathring\mu,\mathring\Lambda)$ is stationarity in $(\mathring\mu,S)$. The covariance condition holds for every $\Delta\in\mathring\Xi$ exactly when $\operatorname{pr}_{\mathring\Xi}(\mathring\Sigma^{-1}-\mathring\Lambda+\mathbb E_R\nabla^2\mathbf U)=0$, that is, since $\mathring\Lambda\in\mathring\Xi$, when $\mathring\Lambda=\operatorname{pr}_{\mathring\Xi}(\mathring\Sigma^{-1}+\mathbb E_R\nabla^2\mathbf U)$. Differentiation under $\mathbb E_{Q_Y}$ gives $\mathbb E_R\nabla\mathbf U=\mathbb E_R\nabla\mathcal A_{\mathrm{loop}}+\mathbb E_{R\otimes Q_Y}\nabla H_o$, and likewise for Hessians.
(c) For fixed $y$ and remaining generators, Lemma~\ref{lem:llp_exp_derivative}(c) gives $\nabla_eH_o=(\mathrm d\rho)^\top\mathbf L((\mathrm d\rho\mathring\omega_e)^\top,\partial H_o/\partial M_e)$. The map $\mathbf S\mapsto(\mathrm d\rho)^\top\mathbf L(\mathbf B,\mathbf S)$ is linear and continuous, so it commutes with $\mathbb E_{Q_Y}$, and the $e$ component of the mean equation becomes Eq.~\eqref{eq:latent_link_current}.
(d) Since $\mathbb E_R(\mathring\omega-\mathring\alpha)=\mathring\mu-\mathring\alpha$, the mean equation is $\mathbb E_R\nabla\mathbf V=0$. Write $\mathring\omega=\mathring\mu+\mathring\Lambda^{-1/2}\mathbf b$. For $\|\mathring\Lambda^{-1}\|\leq1$, $\|\nabla\mathbf V(\mathring\omega)\|$ is dominated by $c(1+\|\mathring\mu\|+\|\mathbf b\|)e^{c(\|\mathring\mu\|+\|\mathbf b\|)}$, so for $\mathring\mu$ in a compact set the contribution of $\{\|\mathbf b\|>c'\}$ is uniformly small for large $c'$, while on $\{\|\mathbf b\|\leq c'\}$ the uniform continuity of $\nabla\mathbf V$ on bounded sets makes $\|\nabla\mathbf V(\mathring\omega)-\nabla\mathbf V(\mathring\mu)\|$ uniformly small; this is the uniform convergence. Along solutions with $\mathring\mu\to\mathring\mu_*$, $\|\nabla\mathbf V(\mathring\mu_*)\|\leq\|\nabla\mathbf V(\mathring\mu_*)-\nabla\mathbf V(\mathring\mu)\|+\|\nabla\mathbf V(\mathring\mu)-\mathbb E_R\nabla\mathbf V\|\to0$. For the last claim, $P_{\mathrm{gen}}\leq Z_{\mathrm{loop}}^{-1}\nu^{\mathcal T}$ gives $D_{\mathrm{KL}}(R\|P_{\mathrm{gen}})\geq D_{\mathrm{KL}}(R\|\nu^{\mathcal T})+\log Z_{\mathrm{loop}}$, so every $R$ with $D_{\mathrm{KL}}(R\|\nu^{\mathcal T})=+\infty$ has infinite free energy by Eq.~\eqref{eq:llp_decomposition}. For the others, Eq.~\eqref{eq:llp_decomposition} and the argument of (a) give Eq.~\eqref{eq:llp_field_free_energy} in $(-\infty,+\infty]$, because $\mathcal A_{\mathrm{loop}}\geq0$ and the matter integrand is bounded below. Since $\mathbf U$ is finite and bounded below, $0<\int e^{-\mathbf U}d\nu^{\mathcal T}<\infty$ and $D_{\mathrm{KL}}(R\|\nu^{\mathcal T})+\mathbb E_R\mathbf U=D_{\mathrm{KL}}(R\|R^\star)-\log\int e^{-\mathbf U}d\nu^{\mathcal T}$, which is minimal exactly at $R=R^\star$; the bound $R^\star\leq c\nu^{\mathcal T}$ makes its free energy finite. The density of $\nu^{\mathcal T}$ is proportional to $e^{-\frac12\langle\mathring\omega-\mathring\alpha,\mathring\Sigma^{-1}(\mathring\omega-\mathring\alpha)\rangle}$.
\end{proof}

\begin{corollary}[The current of a comparison record]\label{cor:llp_delta_rule}
In the setting of Proposition~\ref{prop:llp_field_equation}(c), write samples and links in the reference gauge, in which every frame is the identity; a passive change alters no generator. Let the belief marginals of $Q_Y$ have finite second moments, with means $\mu_i,\mu_j$, covariance $\Sigma_j$ of $k_j$ and cross covariance $\Sigma_{ij}=\mathbb E_{Q_Y}[(k_i-\mu_i)(k_j-\mu_j)^\top]$.
(a) If the only term of $H_o$ that reads $\mathring\omega_e$, $e:j\to i$, is $\frac12(k_i-M_ek_j)^\top W_a(k_i-M_ek_j)$ with $W_a\succ0$, then $\mathbf J_e$ is given by Eq.~\eqref{eq:latent_link_delta_current}.
(b) Consider the selective record~\eqref{eq:selective_record} of row $i$ at the zero comparison record, with the source label $J_i$ among the matter coordinates and the sending map of the source $j$ reached through $e$ equal to $B_{ij}=B_i^{\mathrm{ref}}M_e$. If $Q_Y$ makes $J_i$ independent of the remaining matter coordinates with a row $\beta_i$ that does not depend on $(\mathring\omega,y)$, and the only terms of $H_o$ that read $\mathring\omega_e$ are those of row $i$ with $J_i=j$, then
\begin{equation}
\mathbf J_e=\frac{\beta_{ij}}{\tau_i}(B_i^{\mathrm{ref}})^\top\Gamma_i^{-1}\bigl[(A_i\mu_i-B_i^{\mathrm{ref}}M_e\mu_j)\mu_j^\top+A_i\Sigma_{ij}-B_i^{\mathrm{ref}}M_e\Sigma_j\bigr].
\label{eq:llp_selective_current}
\end{equation}
In both cases Eq.~\eqref{eq:latent_link_current} holds with this current.
\end{corollary}

\begin{proof}
For a matrix $\mathbf M$ in place of $M_e$, the trace-pairing gradient of $\frac12(k_i-\mathbf Mk_j)^\top W_a(k_i-\mathbf Mk_j)$ in $\mathbf M$ is $-W_a(k_i-\mathbf Mk_j)k_j^\top$, which is $Q_Y$-integrable by the second moments, and $\mathbb E_{Q_Y}[(k_i-\mathbf Mk_j)k_j^\top]=\mu_i\mu_j^\top+\Sigma_{ij}-\mathbf M(\Sigma_j+\mu_j\mu_j^\top)$, which gives (a). In (b), by Eq.~\eqref{eq:normalized_link_energy} the row-$i$ part of $H_o$ is $D_{iJ_i}(y)/\tau_i-\log c_i-\log\pi_{iJ_i}$, and $c_i$ does not depend on the links because $\Gamma_i$ is common to the candidate sources. Only $D_{ij}$ reads $M_e$, with gradient $-(B_i^{\mathrm{ref}})^\top\Gamma_i^{-1}(A_ik_i-B_i^{\mathrm{ref}}\mathbf Mk_j)k_j^\top$, so $\partial H_o/\partial M_e$ is $\mathbf 1\{J_i=j\}/\tau_i$ times this matrix; the independence of $J_i$ gives the factor $\beta_{ij}$, and the expectation is computed as in (a). The last sentence is Proposition~\ref{prop:llp_field_equation}(c).
\end{proof}

\subsection{Fisher recognition flow and the learning step}\label{app:llp_flow}

In this subsection $\mathsf G=g^F\oplus\mathsf G_Y$ is the product metric of Section~\ref{sec:latent_link_recognition}. Let the agents' recognition coordinates $\lambda$ range over a finite-dimensional $C^2$ manifold $\mathcal Q_Y$ with a $C^1$ positive-definite metric $\mathsf G_Y$, let $R=\mathcal N(\mathring\mu,\mathring\Lambda^{-1})$ with $\mathring\Lambda$ positive definite in the space $\mathring\Xi$ of Section~\ref{app:llp_field}, and assume $\ell_o\leq c_{\mathfrak A}$ and $p(o\mid\mathring\omega)>0$ for every $\mathring\omega$. With $\mathfrak i_\lambda$ the generator integrand of Eq.~\eqref{eq:latent_link_integrand}, taken without an additive constant, Proposition~\ref{prop:llp_decomposition}(c) with $Q_{\mathring\omega}=Q_\lambda$, the Gaussian density of $\nu^{\mathcal T}$ and the entropy of $R$ give
\begin{equation}
F(\mathring\mu,\mathring\Lambda,\lambda)=\mathcal F_V[R\otimes Q_\lambda;X,o]=\mathbb E_R\mathfrak i_\lambda+\tfrac12\log\det\mathring\Lambda+c_{\mathrm{gen}},
\label{eq:llp_flow_free_energy}
\end{equation}
where $c_{\mathrm{gen}}=\log Z_{\mathrm{loop}}+\frac12\log\det(2\pi\mathring\Sigma)-\frac12\dim\mathcal T\log(2\pi e)$ depends only on the fixed prior. Assume that $F$ is $C^2$ on $\mathcal T\times\mathring\Xi_{\succ0}\times\mathcal Q_Y$, where $\mathring\Xi_{\succ0}$ is the positive-definite cone of $\mathring\Xi$, and that each $\mathfrak i_\lambda$ is $C^2$ with $|\mathfrak i_\lambda|+\|\nabla\mathfrak i_\lambda\|+\|\nabla^2\mathfrak i_\lambda\|\leq c_1e^{c_2\|\mathring\omega\|}$. For $\mathbf a\in\mathcal T$ let $\mathbf a\otimes\mathbf a$ be the operator $\mathring\omega\mapsto\langle\mathbf a,\mathring\omega\rangle\mathbf a$.

\begin{proposition}[Natural-gradient recognition flow of the generators]\label{prop:llp_flow}
(a) The laws $R$ form a minimal regular exponential family with natural parameter $\mathring\vartheta=(\mathring\Lambda\mathring\mu,-\frac12\mathring\Lambda)$, sufficient statistic $(\mathring\omega,\operatorname{pr}_{\mathring\Xi}(\mathring\omega\otimes\mathring\omega))$, expectation parameter $\mathring\varpi=(\mathring\mu,\mathring\Lambda^{-1}+\operatorname{pr}_{\mathring\Xi}(\mathring\mu\otimes\mathring\mu))$ and log-partition function $\mathring\Gamma$. Its Fisher metric is smooth and positive definite, and in the coordinates $(\mathring\mu,\mathring\Lambda)$
\begin{equation}
g^F\bigl((\delta\mathring\mu,\delta\mathring\Lambda),(\delta\mathring\mu,\delta\mathring\Lambda)\bigr)=\langle\delta\mathring\mu,\mathring\Lambda\delta\mathring\mu\rangle+\tfrac12\operatorname{tr}(\mathring\Lambda^{-1}\delta\mathring\Lambda\mathring\Lambda^{-1}\delta\mathring\Lambda).
\label{eq:llp_gaussian_fisher}
\end{equation}
(b) The field $-\operatorname{grad}^{\mathsf G}F$, defined by $\mathsf G(\operatorname{grad}^{\mathsf G}F,Z)=dF[Z]$ for every tangent vector $Z$, is $C^1$; its maximal integral curves exist, are unique, and satisfy $\frac{d}{dt}F=-\|\operatorname{grad}^{\mathsf G}F\|_{\mathsf G}^2$.
(c) The flow is of first order: its velocity is a function of the configuration alone and vanishes exactly at the critical points of $F$. It obeys the force balance $\mathsf G(\dot Q,\cdot)=-dF$, in which $\mathsf G$ acts as a friction tensor. $F$ strictly decreases along every nonconstant solution, so no nonconstant solution is periodic.
(d) In the chart $(\mathring\vartheta,\lambda)$ the generator block of the explicit Euler step with step $\Delta t>0$ is Eq.~\eqref{eq:latent_link_blr}, whenever $\mathring\Lambda'$ is positive definite, as when $0<\Delta t<1$ and $\operatorname{pr}_{\mathring\Xi}\mathbb E_R\nabla^2\mathfrak i_\lambda\succeq0$.
(e) Let $R'=\mathcal N(\mathring\mu',\mathring\Lambda'^{-1})$ have expectation parameter $\mathring\varpi'$, let $\mathring\Gamma^\star$ be the convex conjugate of $\mathring\Gamma$, and let $\widetilde F$ be $F$ at fixed $\lambda$ as a function of $\mathring\varpi$. If $\mathring\Lambda'\succ0$ and $\nabla^2\widetilde F\preceq c_{\mathsf G}\nabla^2\mathring\Gamma^\star$ on the segment from $\mathring\varpi$ to $\mathring\varpi'$, then
\begin{equation}
F(R')\leq F(R)-\Delta t^{-1}D_{\mathrm{KL}}(R\|R')-(\Delta t^{-1}-c_{\mathsf G})D_{\mathrm{KL}}(R'\|R),
\label{eq:llp_blr_descent}
\end{equation}
so a step with $0<\Delta t\leq1/c_{\mathsf G}$ decreases $F$ unless $R'=R$. For the full family and $\mathfrak i_\lambda(\mathring\omega)=\frac12\langle\mathring\omega,\mathbf Q\mathring\omega\rangle-\langle\mathbf b,\mathring\omega\rangle+c_3$ with $\mathbf Q\succ0$, the hypothesis holds with $c_{\mathsf G}=1$, and the step $\Delta t=1$ gives $R'=\mathcal N(\mathbf Q^{-1}\mathbf b,\mathbf Q^{-1})$, the minimizer of $F$ over Gaussian $R$.
\end{proposition}

\begin{proof}
(a) The Lebesgue density of $R$ is $\exp(\langle\mathring\vartheta_1,\mathring\omega\rangle+\langle\mathring\omega,\mathring\vartheta_2\mathring\omega\rangle-\mathring\Gamma(\mathring\vartheta))$ with $\mathring\Gamma(\mathring\vartheta)=-\frac14\langle\mathring\vartheta_1,\mathring\vartheta_2^{-1}\mathring\vartheta_1\rangle-\frac12\log\det(-2\mathring\vartheta_2)+\frac12\dim\mathcal T\log(2\pi)$. For $\mathring\vartheta_2\in\mathring\Xi$, $\langle\mathring\omega,\mathring\vartheta_2\mathring\omega\rangle=\operatorname{tr}(\mathring\vartheta_2(\mathring\omega\otimes\mathring\omega))=\operatorname{tr}(\mathring\vartheta_2\operatorname{pr}_{\mathring\Xi}(\mathring\omega\otimes\mathring\omega))$, which identifies the sufficient statistic. The defining integral is finite exactly when $\mathring\vartheta_2$ is negative definite, an open condition, so the family is regular, and it is minimal because a nonzero polynomial of degree at most two vanishes only on a Lebesgue-null set. Since $\mathring\Lambda^{-1}\in\mathring\Xi$, $\nabla\mathring\Gamma$, the mean of the sufficient statistic, is $\mathring\varpi$~\cite{Brown1986}. The score of $R$ in direction $(\delta\mathring\mu,\delta\mathring\Lambda)$ is $\langle\delta\mathring\mu,\mathring\Lambda(\mathring\omega-\mathring\mu)\rangle-\frac12(\langle\mathring\omega-\mathring\mu,\delta\mathring\Lambda(\mathring\omega-\mathring\mu)\rangle-\operatorname{tr}(\mathring\Lambda^{-1}\delta\mathring\Lambda))$. Its two parts are uncorrelated because odd central moments vanish; the first has variance $\langle\delta\mathring\mu,\mathring\Lambda\delta\mathring\mu\rangle$, and by Isserlis' identity $\langle\mathring\omega-\mathring\mu,\Delta(\mathring\omega-\mathring\mu)\rangle$ has variance $2\operatorname{tr}(\Delta\mathring\Lambda^{-1}\Delta\mathring\Lambda^{-1})$ for self-adjoint $\Delta$, which gives Eq.~\eqref{eq:llp_gaussian_fisher}. The form vanishes only when $\delta\mathring\mu=0$ and $\mathring\Lambda^{-1/2}\delta\mathring\Lambda\mathring\Lambda^{-1/2}=0$, and its coefficients are smooth.
(b) $\mathsf G$ is $C^1$ and positive definite and $F$ is $C^2$, so $\operatorname{grad}^{\mathsf G}F$ is $C^1$, and the Picard--Lindel\"of theorem gives existence and uniqueness of maximal integral curves. Along one, $\frac{d}{dt}F=dF[\dot Q]=-\mathsf G(\operatorname{grad}^{\mathsf G}F,\operatorname{grad}^{\mathsf G}F)$.
(c) The field depends on the configuration only, and its zeros are those of $dF$. The defining relation with $\dot Q=-\operatorname{grad}^{\mathsf G}F$ is the force balance. If the gradient vanished at some time, the constant curve through that point would be a solution, and uniqueness would make the solution constant. On a nonconstant solution the gradient therefore never vanishes and $F$ strictly decreases, so the solution cannot return to a point it has left.
(d) In the natural chart the Fisher metric is $\nabla^2\mathring\Gamma$~\cite{Brown1986}, and the chain rule through $\mathring\varpi=\nabla\mathring\Gamma(\mathring\vartheta)$ gives $\nabla_{\mathring\vartheta}F=\nabla^2\mathring\Gamma\nabla_{\mathring\varpi}F$. Since $\mathsf G$ is block diagonal, the generator block of the flow is $\dot{\mathring\vartheta}=-\nabla_{\mathring\varpi}F$, with Euler step $\mathring\vartheta'=\mathring\vartheta-\Delta t\nabla_{\mathring\varpi}F$. By Lemma~\ref{lem:llp_gaussian_gradients}, $\nabla_{\mathring\mu}\mathbb E_R\mathfrak i_\lambda=\mathbb E_R\nabla\mathfrak i_\lambda$, and the gradient of $\mathbb E_R\mathfrak i_\lambda$ in the covariance $S=\mathring\Lambda^{-1}$ within $\mathring\Xi$ is $\frac12\operatorname{pr}_{\mathring\Xi}\mathbb E_R\nabla^2\mathfrak i_\lambda$, while that of $\frac12\log\det\mathring\Lambda=-\frac12\log\det S$ is $-\frac12\mathring\Lambda$. Put $\Upsilon=\frac12\operatorname{pr}_{\mathring\Xi}\mathbb E_R\nabla^2\mathfrak i_\lambda-\frac12\mathring\Lambda\in\mathring\Xi$. Since $\mathring\mu=\mathring\varpi_1$ and $S=\mathring\varpi_2-\operatorname{pr}_{\mathring\Xi}(\mathring\varpi_1\otimes\mathring\varpi_1)$, and since $\operatorname{tr}(\Upsilon\operatorname{pr}_{\mathring\Xi}\Delta)=\operatorname{tr}(\Upsilon\Delta)$ and $\operatorname{tr}(\Upsilon(\eta\otimes\mathring\mu+\mathring\mu\otimes\eta))=2\langle\Upsilon\mathring\mu,\eta\rangle$, the chain rule gives $\nabla_{\mathring\varpi_2}F=\Upsilon$ and $\nabla_{\mathring\varpi_1}F=\mathbb E_R\nabla\mathfrak i_\lambda-2\Upsilon\mathring\mu$. The second component of the Euler step is $-\frac12\mathring\Lambda'=-\frac12\mathring\Lambda-\Delta t\Upsilon$, the first equation of Eq.~\eqref{eq:latent_link_blr}. The first component is $\mathring\Lambda'\mathring\mu'=\mathring\Lambda\mathring\mu-\Delta t\mathbb E_R\nabla\mathfrak i_\lambda+2\Delta t\Upsilon\mathring\mu=\mathring\Lambda'\mathring\mu-\Delta t\mathbb E_R\nabla\mathfrak i_\lambda$, the second. If $0<\Delta t<1$ and the projected expected Hessian is positive semidefinite, $\mathring\Lambda'\succeq(1-\Delta t)\mathring\Lambda\succ0$.
(e) The Lebesgue density $r$ of $R$ gives $\mathbb E_R\log r=\langle\mathring\vartheta,\mathring\varpi\rangle-\mathring\Gamma(\mathring\vartheta)=\mathring\Gamma^\star(\mathring\varpi)$~\cite{Rockafellar1970}, so Eq.~\eqref{eq:llp_flow_free_energy} reads $\widetilde F=\mathbb E_R\mathfrak i_\lambda+\mathring\Gamma^\star+\mathrm{const}$. Taking the $R$-expectation of $\log(r/r')$ gives $D_{\mathrm{KL}}(R\|R')=\mathring\Gamma(\mathring\vartheta')-\mathring\Gamma(\mathring\vartheta)-\langle\mathring\varpi,\mathring\vartheta'-\mathring\vartheta\rangle$, and the Bregman divergence of $\mathring\Gamma^\star$ is $\mathring\Gamma^\star(\mathring\varpi')-\mathring\Gamma^\star(\mathring\varpi)-\langle\mathring\vartheta,\mathring\varpi'-\mathring\varpi\rangle=D_{\mathrm{KL}}(R'\|R)$. The expectation domain $\{\mathring\varpi_2-\operatorname{pr}_{\mathring\Xi}(\mathring\varpi_1\otimes\mathring\varpi_1)\succ0\}$ is convex: along a segment with weights $1-s$ and $s$, this difference equals the convex combination of its endpoint values plus $s(1-s)\operatorname{pr}_{\mathring\Xi}((\mathbf a-\mathbf a')\otimes(\mathbf a-\mathbf a'))$, where $\mathbf a,\mathbf a'$ are the endpoint means, and $\operatorname{pr}_{\mathring\Xi}$ preserves positive semidefiniteness. Taylor's formula with integral remainder along the segment and the curvature hypothesis give $\widetilde F(\mathring\varpi')\leq\widetilde F(\mathring\varpi)+\langle\nabla\widetilde F(\mathring\varpi),\mathring\varpi'-\mathring\varpi\rangle+c_{\mathsf G}D_{\mathrm{KL}}(R'\|R)$, since the same formula applied to $\mathring\Gamma^\star$ produces its Bregman divergence. By (d), $\nabla\widetilde F(\mathring\varpi)=\Delta t^{-1}(\mathring\vartheta-\mathring\vartheta')$, and adding the two divergence formulas gives $\langle\mathring\vartheta'-\mathring\vartheta,\mathring\varpi'-\mathring\varpi\rangle=D_{\mathrm{KL}}(R\|R')+D_{\mathrm{KL}}(R'\|R)$; substitution proves Eq.~\eqref{eq:llp_blr_descent}. For the quadratic integrand in the full family, $\mathbb E_R\mathfrak i_\lambda=\frac12\operatorname{tr}(\mathbf Q\mathring\varpi_2)-\langle\mathbf b,\mathring\varpi_1\rangle+c_3$ is affine in $\mathring\varpi$, so $\nabla^2\widetilde F=\nabla^2\mathring\Gamma^\star$. Moreover $\mathbb E_R\nabla^2\mathfrak i_\lambda=\mathbf Q$ and $\mathbb E_R\nabla\mathfrak i_\lambda=\mathbf Q\mathring\mu-\mathbf b$, so Eq.~\eqref{eq:latent_link_blr} with $\Delta t=1$ gives $\mathring\Lambda'=\mathbf Q$ and $\mathring\mu'=\mathbf Q^{-1}\mathbf b$. Then $R'\propto e^{-\mathfrak i_\lambda}$, and $F=D_{\mathrm{KL}}(R\|R')+\mathrm{const}$ is minimal at $R=R'$.
\end{proof}

\subsection{The episode-history prior and history inertia}\label{app:llp_history}

Use the history law~\eqref{eq:latent_link_history_law} with $\mathcal A_{\mathrm{loop}}$ finite, measurable and nonnegative, $\mathbf m>0$, and either kinetic term: the Gaussian increment $\mathfrak U(\mathring\omega',\mathring\omega)=\frac12\|\mathring\omega'-\mathring\omega\|^2$, or the time-like defect $\mathfrak U(\mathring\omega',\mathring\omega)=\sum_{e\in E_+}\mathfrak Q(e^{-\mathrm d\rho\mathring\omega'_e}e^{\mathrm d\rho\mathring\omega_e})$ with $\mathfrak Q(\mathbf H)=\frac12\|\mathbf H-I\|_{\mathsf G}^2$. The matrix $e^{-\mathrm d\rho\mathring\omega'_e}e^{\mathrm d\rho\mathring\omega_e}$ is the represented reference holonomy of the time-like loop formed by $e$ at one episode, $\bar e$ at the next and the identity links of its endpoints between them.

\begin{proposition}[The generator history prior]\label{prop:llp_history_prior}
(a) $0<Z_{\mathrm{hist}}\leq Z_{\mathrm{loop}}^{n_1+1}\leq1$, so $\mathbb P^{\mathrm{hist}}$ is a probability law with density at most $Z_{\mathrm{hist}}^{-1}$ with respect to $\bigotimes_n\nu^{\mathcal T}$, and every function bounded by $c_1e^{c_2\|\mathring\omega^{(0:n_1)}\|}$ is integrable under it.
(b) Put $Z^{\leftarrow}_{n_1}\equiv1$ and, for $n<n_1$, $Z^{\leftarrow}_n(\mathring\omega)=\int e^{-\mathcal A_{\mathrm{loop}}(\mathring\omega')-\mathbf m\mathfrak U(\mathring\omega',\mathring\omega)}Z^{\leftarrow}_{n+1}(\mathring\omega')\nu^{\mathcal T}(d\mathring\omega')$. Then $\mathbb P^{\mathrm{hist}}$ is the law of the Markov chain with initial law $Z_{\mathrm{hist}}^{-1}e^{-\mathcal A_{\mathrm{loop}}}Z^{\leftarrow}_0\nu^{\mathcal T}$ and transition kernels
\begin{equation}
\mathbb P^{(n+1\leftarrow n)}(d\mathring\omega'\mid\mathring\omega)=\frac{e^{-\mathcal A_{\mathrm{loop}}(\mathring\omega')-\mathbf m\mathfrak U(\mathring\omega',\mathring\omega)}Z^{\leftarrow}_{n+1}(\mathring\omega')}{Z^{\leftarrow}_n(\mathring\omega)}\nu^{\mathcal T}(d\mathring\omega'),
\label{eq:llp_history_kernels}
\end{equation}
which depend on $n$ and on the horizon through the backward normalizers. The chain whose kernels normalize each increment factor separately, started from the law proportional to $e^{-\mathcal A_{\mathrm{loop}}}\nu^{\mathcal T}$, has the law $\mathbb P^{\mathrm{hist}}$ exactly when $Z^{\leftarrow}_{n_1-1}$ is constant $\nu^{\mathcal T}$-almost everywhere; this fails for one scalar generator with $\nu^{\mathcal T}=\mathcal N(0,1)$, $\mathcal A_{\mathrm{loop}}=0$, the Gaussian increment and $n_1=1$.
(c) With static frames, the identity link of agent $i$ between two episodes has the constant twist $U_i^{-1}U_i=I$. The time-like links therefore carry no variable, $\mathbb P^{\mathrm{hist}}$ is a law on the spatial generator histories alone, and each kinetic term is a function of the generators of one edge at two consecutive episodes.
\end{proposition}

\begin{proof}
(a) Since $\mathcal A_{\mathrm{loop}}\geq0$, $\mathfrak U\geq0$ and $\mathbf m>0$, the integrand of $Z_{\mathrm{hist}}$ is positive and at most $\prod_ne^{-\mathcal A_{\mathrm{loop}}(\mathring\omega^{(n)})}$, whose integral is $Z_{\mathrm{loop}}^{n_1+1}\leq1$; a product of Gaussian laws integrates $e^{c_2\|\cdot\|}$.
(b) By Tonelli's theorem and induction from $n_1$, each $Z^{\leftarrow}_n$ is measurable with values in $(0,1]$, each kernel has total mass one, and $\int e^{-\mathcal A_{\mathrm{loop}}}Z^{\leftarrow}_0d\nu^{\mathcal T}=Z_{\mathrm{hist}}$. Along a history the ratios $Z^{\leftarrow}_{n+1}/Z^{\leftarrow}_n$ telescope and $Z^{\leftarrow}_{n_1}\equiv1$, so the chain density is that of $\mathbb P^{\mathrm{hist}}$, and the kernels are its conditional laws. The kernels of the separately normalized chain are the increment factors divided by $Z^{\leftarrow}_{n_1-1}(\mathring\omega)$, the normalizer of one factor, so that chain has density proportional to $\prod_{n<n_1}Z^{\leftarrow}_{n_1-1}(\mathring\omega^{(n)})^{-1}$ relative to $\mathbb P^{\mathrm{hist}}$; by Fubini's theorem this product is almost everywhere constant exactly when $Z^{\leftarrow}_{n_1-1}$ is. In the scalar case $Z^{\leftarrow}_0(\mathring\omega)=(1+\mathbf m)^{-1/2}\exp[-\mathbf m\mathring\omega^2/(2(1+\mathbf m))]$, which is not constant.
(c) With a static frame $U_i$, the twist of the identity link does not depend on any variable, and Eq.~\eqref{eq:latent_link_history_law} contains no other link.
\end{proof}

\begin{proposition}[Inertia of the most probable history]\label{prop:llp_inertia}
Let $\mathfrak V_n:\mathcal T\to\mathbb R$ be differentiable and $\mathcal A^{\mathrm{hist}}=\sum_{n=0}^{n_1}\mathfrak V_n(\mathring\omega^{(n)})+\mathbf m\sum_{n<n_1}\mathfrak U(\mathring\omega^{(n+1)},\mathring\omega^{(n)})$. With $\mathfrak V_n=\mathcal A_{\mathrm{loop}}+\frac12\langle\cdot-\mathring\alpha,\mathring\Sigma^{-1}(\cdot-\mathring\alpha)\rangle$, $\mathcal A^{\mathrm{hist}}$ is the negative log Lebesgue density of $\mathbb P^{\mathrm{hist}}$ up to a constant, and adding an episode matter term gives the corresponding posterior or conditional objective. A most probable history minimizes $\mathcal A^{\mathrm{hist}}$ with free or held ends.
(a) For the Gaussian increment every critical point satisfies Eq.~\eqref{eq:latent_link_history_inertia}, and with free ends $\mathbf m(\mathring\omega^{(0)}-\mathring\omega^{(1)})+\nabla\mathfrak V_0(\mathring\omega^{(0)})=0$ and $\mathbf m(\mathring\omega^{(n_1)}-\mathring\omega^{(n_1-1)})+\nabla\mathfrak V_{n_1}(\mathring\omega^{(n_1)})=0$. If each $\mathfrak V_n$ is continuous and $\sum_n\mathfrak V_n(\mathring\omega^{(n)})\to\infty$ as $\|\mathring\omega^{(0:n_1)}\|\to\infty$, as when each $\mathfrak V_n$ is at least the anchor quadratic minus a constant, a most probable history exists.
(b) For the time-like defect let $\mathrm d\rho$ be injective, let $\mathfrak g$ carry the pulled-back inner product, let $\operatorname{pr}_{\mathfrak g}$ be the $\langle\cdot,\cdot\rangle_{\mathsf G}$-orthogonal projection of matrices onto $\mathrm d\rho(\mathfrak g)$, read in $\mathfrak g$, and write $\Theta^{\leftarrow}_e=e^{-\mathrm d\rho\mathring\omega'_e}e^{\mathrm d\rho\mathring\omega_e}$ for a later family $\mathring\omega'$ and an earlier family $\mathring\omega$. Then
\begin{equation}
\begin{aligned}
(\partial_{\mathring\omega}\mathfrak U)_e&=\operatorname{pr}_{\mathfrak g}\mathbf L\bigl((\mathrm d\rho\mathring\omega_e)^*,(e^{-\mathrm d\rho\mathring\omega'_e})^*(\Theta^{\leftarrow}_e-I)\bigr),\\
(\partial_{\mathring\omega'}\mathfrak U)_e&=-\operatorname{pr}_{\mathfrak g}\mathbf L\bigl(-(\mathrm d\rho\mathring\omega'_e)^*,(\Theta^{\leftarrow}_e-I)(e^{\mathrm d\rho\mathring\omega_e})^*\bigr),
\end{aligned}
\label{eq:llp_timelike_gradients}
\end{equation}
and every critical point satisfies, for $0<n<n_1$,
\begin{equation}
\mathbf m\bigl[\partial_{\mathring\omega}\mathfrak U(\mathring\omega^{(n+1)},\mathring\omega^{(n)})+\partial_{\mathring\omega'}\mathfrak U(\mathring\omega^{(n)},\mathring\omega^{(n-1)})\bigr]+\nabla\mathfrak V_n(\mathring\omega^{(n)})=0.
\label{eq:llp_timelike_inertia}
\end{equation}
For $r>0$ let $c_r=\sqrt2\sup\|D^3\mathfrak U\|$ over $\|\mathring\omega\|,\|\mathring\omega'\|\leq r$. On that set $\|\partial_{\mathring\omega}\mathfrak U-(\mathring\omega-\mathring\omega')\|$ and $\|\partial_{\mathring\omega'}\mathfrak U-(\mathring\omega'-\mathring\omega)\|$ are at most $c_rr\|\mathring\omega'-\mathring\omega\|$, so at a history in the ball of radius $r$ the left side of Eq.~\eqref{eq:llp_timelike_inertia} differs from that of Eq.~\eqref{eq:latent_link_history_inertia} by at most $\mathbf mc_rr(\|\mathring\omega^{(n+1)}-\mathring\omega^{(n)}\|+\|\mathring\omega^{(n)}-\mathring\omega^{(n-1)}\|)$. At fixed $\mathring\omega$, $\mathfrak U(\mathring\omega+\delta\mathring\omega,\mathring\omega)=\frac12\sum_e\|\mathbf L(-\mathrm d\rho\mathring\omega_e,\mathrm d\rho\delta\mathring\omega_e)e^{\mathrm d\rho\mathring\omega_e}\|_{\mathsf G}^2+O(\|\delta\mathring\omega\|^3)$, a generator-dependent increment tensor that reduces to $\frac12\|\delta\mathring\omega\|^2$ at $\mathring\omega=0$.
\end{proposition}

\begin{proof}
(a) $\mathring\omega^{(n)}$ enters the increments to $n-1$ and $n+1$ in the interior and one increment at each free end, so differentiation gives the equations. Under the growth condition $\mathcal A^{\mathrm{hist}}$ is continuous and coercive on a finite-dimensional space, or on the affine subspace of held ends, and attains its minimum.
(b) By Lemma~\ref{lem:llp_exp_derivative}(a), the differential of $\Theta^{\leftarrow}_e$ in direction $\delta\mathring\omega_e$ is $e^{-\mathrm d\rho\mathring\omega'_e}\mathbf L(\mathrm d\rho\mathring\omega_e,\mathrm d\rho\delta\mathring\omega_e)$ in the earlier argument and $-\mathbf L(-\mathrm d\rho\mathring\omega'_e,\mathrm d\rho\delta\mathring\omega_e)e^{\mathrm d\rho\mathring\omega_e}$ in the later one, and the differential of $\mathfrak Q$ at $\mathbf H$ is $\mathbf T\mapsto\langle\mathbf H-I,\mathbf T\rangle_{\mathsf G}$. Moving the factors across the inner product with Lemmas~\ref{lem:llp_trace_form}(a) and~\ref{lem:llp_exp_derivative}(b), and projecting onto $\mathrm d\rho(\mathfrak g)$, gives Eq.~\eqref{eq:llp_timelike_gradients}, and Eq.~\eqref{eq:llp_timelike_inertia} is the stationarity of $\mathcal A^{\mathrm{hist}}$. The function $\mathfrak U$ is smooth, nonnegative and zero on the diagonal, so its gradient vanishes there. Since $\Theta^{\leftarrow}_e=I+\mathrm d\rho(\mathring\omega_e-\mathring\omega'_e)+O(\|(\mathring\omega,\mathring\omega')\|^2)$ and the inner product is pulled back, $\mathfrak U=\frac12\|\mathring\omega-\mathring\omega'\|^2+O(\|(\mathring\omega,\mathring\omega')\|^3)$, so the mixed second derivative at the origin is $-I$. Let $\mathrm R_1(\mathring\omega',\mathring\omega)=\partial_{\mathring\omega}\mathfrak U-(\mathring\omega-\mathring\omega')$. Then $\mathrm R_1(\mathring\omega,\mathring\omega)=0$ and $\mathrm R_1(\mathring\omega',\mathring\omega)=\int_0^1D_{\mathring\omega'}\mathrm R_1(\mathring\omega+t(\mathring\omega'-\mathring\omega),\mathring\omega)[\mathring\omega'-\mathring\omega]dt$, where $D_{\mathring\omega'}\mathrm R_1$ vanishes at the origin and has derivative bounded by $\sup\|D^3\mathfrak U\|$ on the convex set $\{\|\mathring\omega\|,\|\mathring\omega'\|\leq r\}$, whose points lie within $\sqrt2r$ of the origin. This gives the bound, and the same argument applies to $\partial_{\mathring\omega'}\mathfrak U$. For the increment tensor, the differential in the later argument gives $\Theta^{\leftarrow}_e=I-\mathbf L(-\mathrm d\rho\mathring\omega_e,\mathrm d\rho\delta\mathring\omega_e)e^{\mathrm d\rho\mathring\omega_e}+O(\|\delta\mathring\omega\|^2)$ at $\mathring\omega'=\mathring\omega+\delta\mathring\omega$, and $\mathbf L(0,\mathbf T)=\mathbf T$.
\end{proof}

\subsection{The linearized link equation}\label{app:llp_linear}

Assume the setting of Section~\ref{sec:latent_link_proca}: $\mathrm d\rho$ is injective and $\mathfrak g$ carries the pulled-back inner product; the history prior has zero anchor, prior covariance $\mathring\Sigma$ at every episode, action $\mathcal A_{\mathrm{loop}}=\sum_pc_p\mathfrak Q_p$, mass $\mathbf m$ and either kinetic term; and samples are written in the reference gauge. At episode $n$ a finite set $\mathfrak A_n$ of records enters with generator-independent weights $\beta_a\in[0,1]$. Record $a$ has a scope $\{i,j\}$ joined by an edge copy $e:j\to i$ of $E_+$ or of its reversal, with $\mathring\omega_{\bar e}=-\mathring\omega_e$, and the energy $E_a(k;\mathring\omega)=\frac12(k_i-e^{\mathrm d\rho\mathring\omega_e}k_j)^\top W_a(k_i-e^{\mathrm d\rho\mathring\omega_e}k_j)$ with $W_a\succ0$. Matter recognition is a fixed product Gaussian law under which $k_i$ at episode $n$ has mean $\mu_i^{(n)}$ and covariance $\Sigma_i^{(n)}$, and
\begin{equation}
\begin{aligned}
\mathcal A^{\mathrm{hist}}(\mathring\omega)={}&\sum_{n=0}^{n_1}\Bigl[\tfrac12\langle\mathring\omega^{(n)},\mathring\Sigma^{-1}\mathring\omega^{(n)}\rangle+\mathcal A_{\mathrm{loop}}(\mathring\omega^{(n)})+\sum_{a\in\mathfrak A_n}\beta_a\mathbb E[E_a(k;\mathring\omega^{(n)})]\Bigr]\\
&+\mathbf m\sum_{n<n_1}\mathfrak U(\mathring\omega^{(n+1)},\mathring\omega^{(n)}).
\end{aligned}
\label{eq:llp_history_objective}
\end{equation}
On $\mathcal T^{n_1+1}$, with $\langle\mathring\omega,\mathring\omega'\rangle=\sum_n\langle\mathring\omega^{(n)},\mathring\omega'^{(n)}\rangle$, let $\mathbf K$, $\mathring\Phi$, $\mathbf j$ and $\mathbf N_0$ be defined by Eqs.~\eqref{eq:latent_link_proca}, \eqref{eq:latent_link_matter_operator} and~\eqref{eq:latent_link_current_correction}, where $B_{a,n}=\mathrm d\rho\mathring\omega^{(n)}_e$ for the copy $e$ of record $a$, and let $c_{\mathring\Sigma}$ be the smallest eigenvalue of $\mathring\Sigma^{-1}$.

\begin{proposition}[The linearized link equation]\label{prop:llp_linearized}
(a) $\nabla\mathcal A^{\mathrm{hist}}(0)=-\mathbf j$ and $\nabla^2\mathcal A^{\mathrm{hist}}(0)=\mathbf K+\mathbf N_0$, so to first order in the generators and the currents the stationarity condition is $\mathbf K\mathring\omega=\mathbf j$.
(b) $\mathbf K\succeq c_{\mathring\Sigma}I$, so $\mathring\omega_{\mathrm{lin}}=\mathbf K^{-1}\mathbf j$ is the unique solution, with $\|\mathring\omega_{\mathrm{lin}}\|\leq\|\mathbf j\|/c_{\mathring\Sigma}$.
(c) Let $c_0=\|\mathbf N_0\|$ and $c_r=\sup_{\|\mathring\omega\|\leq r}\|D^3\mathcal A^{\mathrm{hist}}(\mathring\omega)\|$. If $c_0+c_rr\leq c_{\mathring\Sigma}/2$ and $\|\mathbf j\|\leq c_{\mathring\Sigma}r/4$, then $\mathcal A^{\mathrm{hist}}$ has exactly one critical point $\mathring\omega_\star$ in the ball of radius $r$ about zero; it is the unique minimizer of $\mathcal A^{\mathrm{hist}}$ on that ball, $\|\mathring\omega_\star\|\leq2\|\mathbf j\|/c_{\mathring\Sigma}$, and
\begin{equation}
\|\mathring\omega_\star-\mathring\omega_{\mathrm{lin}}\|\leq\frac{2\|\mathbf j\|}{c_{\mathring\Sigma}^2}\Bigl(c_0+\frac{c_r\|\mathbf j\|}{c_{\mathring\Sigma}}\Bigr).
\label{eq:llp_linearization_error}
\end{equation}
\end{proposition}

\begin{proof}
(a) The anchor term has gradient zero and Hessian $\mathring\Sigma^{-1}$ at zero. Each $\mathfrak Q_p$ is smooth, being a polynomial in the entries of matrix exponentials, and $\mathfrak Q_p=\frac12\|\mathbf C_p\mathring\omega\|^2+O(\|\mathring\omega\|^3)$ by Proposition~\ref{prop:llp_loop_defect}(c), so by uniqueness of Taylor coefficients its gradient at zero vanishes and its Hessian is $\mathbf C_p^\top\mathbf C_p$. The Gaussian increment is the quadratic form of $\mathbf m(-\Delta_n)$, and the time-like defect has the same second-order part and zero gradient at the origin by the expansion in the proof of Proposition~\ref{prop:llp_inertia}(b). For a record $a$ with $\Omega=\mathrm d\rho\mathring\omega_e$, dropping the episode index, independence of $k_i$ and $k_j$ gives $\mathbb E[E_a]=\frac12[\|\mu_i-e^{\Omega}\mu_j\|_{W_a}^2+\operatorname{tr}(W_a\Sigma_i)+\operatorname{tr}(W_ae^\Omega\Sigma_je^{\Omega^\top})]$. With $e^\Omega=I+\Omega+\frac12\Omega^2+O(\|\Omega\|^3)$ and $M_j=\Sigma_j+\mu_j\mu_j^\top$, the first-order part is $\operatorname{tr}(W_a\Omega\Sigma_j)-(\mu_i-\mu_j)^\top W_a\Omega\mu_j=-\operatorname{tr}(\mathbf J_a^\top\Omega)$, and the second-order part is
\begin{equation*}
\tfrac12\bigl[\operatorname{tr}(W_a\Omega M_j\Omega^\top)+\operatorname{tr}(W_a\Omega^2\Sigma_j)-(\mu_i-\mu_j)^\top W_a\Omega^2\mu_j\bigr]=\tfrac12\bigl[\operatorname{tr}(W_a\Omega M_j\Omega^\top)-\operatorname{tr}(\mathbf J_a^\top\Omega^2)\bigr],
\end{equation*}
because $\mathbf J_a^\top=(\mu_j(\mu_i-\mu_j)^\top-\Sigma_j)W_a$. For a record on a reversed copy, $\Omega=-\mathrm d\rho\mathring\omega_{\bar e}$ with $\bar e\in E_+$, which changes the sign of the first-order part only. Summing with the weights gives the gradient and the Hessian.
(b) The loop, episode and matter forms are nonnegative, the last because $\operatorname{tr}(W_a\Omega M_j\Omega^\top)=\mathbb E\|W_a^{1/2}\Omega k_j\|^2$, so $\langle\mathring\omega,\mathbf K\mathring\omega\rangle\geq c_{\mathring\Sigma}\|\mathring\omega\|^2$ and $\|\mathbf K^{-1}\|\leq1/c_{\mathring\Sigma}$.
(c) Write $\nabla\mathcal A^{\mathrm{hist}}(\mathring\omega)=-\mathbf j+\mathbf K\mathring\omega+\nabla\mathrm R(\mathring\omega)$, so that $\nabla\mathrm R(0)=0$, $\nabla^2\mathrm R(0)=\mathbf N_0$, and $\|\nabla^2\mathrm R(\mathring\omega)\|\leq c_0+c_r\|\mathring\omega\|$ for $\|\mathring\omega\|\leq r$. The critical points are the fixed points of $\mathring\omega\mapsto\mathbf K^{-1}(\mathbf j-\nabla\mathrm R(\mathring\omega))$. On the ball of radius $r_1=2\|\mathbf j\|/c_{\mathring\Sigma}\leq r/2$, $\|\nabla\mathrm R(\mathring\omega)\|\leq(c_0+c_rr)\|\mathring\omega\|\leq\frac12c_{\mathring\Sigma}r_1$, so the map sends this ball into itself, its values having norm at most $\|\mathbf j\|/c_{\mathring\Sigma}+r_1/2=r_1$, and it has Lipschitz constant at most $(c_0+c_rr)/c_{\mathring\Sigma}\leq\frac12$ there; it therefore has exactly one fixed point $\mathring\omega_\star$ in that ball. On the ball of radius $r$, $\nabla^2\mathcal A^{\mathrm{hist}}\succeq(c_{\mathring\Sigma}-c_0-c_rr)I\succeq\frac12c_{\mathring\Sigma}I$, so $\mathcal A^{\mathrm{hist}}$ is strongly convex on that convex set and has at most one critical point there, which is its unique minimizer on the set. Finally $\mathring\omega_\star-\mathring\omega_{\mathrm{lin}}=-\mathbf K^{-1}\nabla\mathrm R(\mathring\omega_\star)$ and $\|\nabla\mathrm R(\mathring\omega_\star)\|\leq(c_0+\frac12c_r\|\mathring\omega_\star\|)\|\mathring\omega_\star\|$, which with $\|\mathring\omega_\star\|\leq r_1$ gives Eq.~\eqref{eq:llp_linearization_error}.
\end{proof}

\subsection{Induced link action}\label{app:llp_induced}

For the undirected alternative, let each link family $\Theta$ carry a nonnegative potential $\widetilde p_{X,\Theta}(o,y)$, jointly measurable in $(\Theta,o,y)$, with matter partition function $Z_{X,\Theta}=\int\widetilde p_{X,\Theta}d(\nu^Y\otimes\nu^O)\in[0,+\infty]$. When $0<Z_{X,\Theta}<\infty$ the undirected model at $\Theta$ is $\widetilde p_{X,\Theta}/Z_{X,\Theta}$.

\begin{proposition}[Induced link action]\label{prop:llp_induced}
(a) The generator marginal of the joint~\eqref{eq:pifb_latent_population_joint} is $P_{\mathrm{gen}}$, and with record densities the link marginal of the posterior is proportional to $p(o\mid\mathring\omega)P_{\mathrm{gen}}(d\mathring\omega)$.
(b) Let $\mathfrak Z$ be as in Eq.~\eqref{eq:pifb_joint_link_matter_normalizer} and assume $0<\mathfrak Z<\infty$. The law with density $\mathfrak Z^{-1}e^{-\mathcal A_{\mathrm{loop}}(\mathring\omega)}\widetilde p_{X,\Theta(\mathring\omega)}(o,y)$ with respect to $\nu^{\mathcal T}\otimes\nu^Y\otimes\nu^O$ is a probability law; $Z_{X,\Theta(\mathring\omega)}<\infty$ for $\nu^{\mathcal T}$-almost every $\mathring\omega$; the generator marginal is $\mathfrak Z^{-1}e^{-\mathcal A_{\mathrm{loop}}-\mathcal A_{\mathrm{ind}}}\nu^{\mathcal T}$ with $\mathcal A_{\mathrm{ind}}=-\log Z_{X,\Theta(\cdot)}\in(-\infty,+\infty]$ almost everywhere; and on $\{0<Z_{X,\Theta(\mathring\omega)}<\infty\}$ the conditional law given $\mathring\omega$ is the undirected model at $\Theta(\mathring\omega)$. A record extension with densities is the case $Z_{X,\Theta}\equiv1$, in which $\mathcal A_{\mathrm{ind}}=0$.
(c) A proper generator prior does not imply $\mathfrak Z<\infty$. For $G=\mathrm{GL}^+(1)$, trivial frames, $\mathcal A_{\mathrm{loop}}=0$, $\nu^{\mathcal T}=\mathcal N(0,1)$, two scalar agents with standard Gaussian self potentials, no record, and the edge potential $\exp(e^{\mathring\omega_e}k_1k_2)$, one has $Z_{X,\Theta}=(1-e^{2\mathring\omega_e})^{-1/2}$ for $\mathring\omega_e<0$ and $Z_{X,\Theta}=+\infty$ for $\mathring\omega_e\geq0$, so $\mathfrak Z=+\infty$.
(d) If the potential is a centered Gaussian, $\widetilde p_{X,\Theta}(y)=\exp(-\frac12y^\top\mathbf P(\Theta)y)$ on $\mathbb R^d$ with symmetric $\mathbf P(\Theta)$, Lebesgue reference and no record, then $Z_{X,\Theta}=(2\pi)^{d/2}(\det\mathbf P(\Theta))^{-1/2}$ when $\mathbf P(\Theta)\succ0$ and $Z_{X,\Theta}=+\infty$ otherwise. Hence $\mathcal A_{\mathrm{ind}}=\frac12\log\det\mathbf P(\Theta(\mathring\omega))-\frac d2\log(2\pi)$, and $\mathfrak Z<\infty$ requires $\mathbf P(\Theta(\mathring\omega))\succ0$ for $\nu^{\mathcal T}$-almost every $\mathring\omega$ and $\int e^{-\mathcal A_{\mathrm{loop}}}(\det\mathbf P(\Theta(\cdot)))^{-1/2}d\nu^{\mathcal T}<\infty$.
\end{proposition}

\begin{proof}
(a) By Tonelli's theorem, integrate the records, each a kernel of unit mass, and then the matter against the probability kernel $P^Y_{\mathring\omega}$; the posterior marginal is that of Proposition~\ref{prop:llp_decomposition}(c).
(b) Tonelli's theorem gives $\int e^{-\mathcal A_{\mathrm{loop}}}\widetilde p_{X,\Theta(\mathring\omega)}d(\nu^Y\otimes\nu^O)=e^{-\mathcal A_{\mathrm{loop}}(\mathring\omega)}Z_{X,\Theta(\mathring\omega)}$, whose $\nu^{\mathcal T}$-integral is $\mathfrak Z$, so the density integrates to one, and finiteness of $\mathfrak Z$ with $e^{-\mathcal A_{\mathrm{loop}}}>0$ forces $Z_{X,\Theta(\mathring\omega)}<\infty$ almost everywhere. The generator marginal is $\mathfrak Z^{-1}e^{-\mathcal A_{\mathrm{loop}}}Z_{X,\Theta(\mathring\omega)}\nu^{\mathcal T}$, with $\mathcal A_{\mathrm{ind}}=+\infty$ where $Z_{X,\Theta}=0$, and dividing the joint density by the marginal density gives the conditional law. The densities of a record extension are conditional densities of unit mass in $(y,o)$ for every link family, so as potentials they have $Z_{X,\Theta}=1$.
(c) Put $s=e^{\mathring\omega_e}$ and $u=(k_1+k_2)/\sqrt2$, $v=(k_1-k_2)/\sqrt2$. The exponent $-(k_1^2+k_2^2)/2+sk_1k_2$ equals $-(1-s)u^2/2-(1+s)v^2/2$, so the integral of $(2\pi)^{-1}$ times its exponential is $((1-s)(1+s))^{-1/2}$ for $s<1$ and diverges in $u$ for $s\geq1$; and $\nu^{\mathcal T}\{\mathring\omega_e\geq0\}=\frac12$.
(d) Diagonalizing $\mathbf P(\Theta)$ factorizes the integral into one-dimensional Gaussian integrals, each equal to $(2\pi/\lambda_k)^{1/2}$ when the eigenvalue $\lambda_k$ is positive and infinite otherwise. The conditions on $\mathfrak Z$ are those of (b).
\end{proof}

\subsection{Coarse link laws}\label{app:llp_coarse}

Let $G$ be a closed subgroup of $\mathrm{GL}(K,\mathbb R)$ acting by inclusion, so that $\mathfrak g\subseteq\mathfrak{gl}(K,\mathbb R)$ and the exponential is the matrix exponential; extend $\langle\cdot,\cdot\rangle_{\mathsf G}$ to $\mathfrak{gl}(K,\mathbb R)$ by the same formula, and let $\|\cdot\|$ be the operator norm and $E_{kl}$ the matrix units. Let the fine generators $\mathring\omega_{e'}$ and $\mathring\omega_e$ of Eq.~\eqref{eq:pifb_coarse_link_product} be independent Gaussian random elements of $\mathfrak g$ with means $\mathring\mu_{e'},\mathring\mu_e$ and positive-definite precisions $\mathring\Lambda_{e'},\mathring\Lambda_e$, self-adjoint for $\langle\cdot,\cdot\rangle_{\mathsf G}$, as under a product Gaussian recognition law of the generators. A coarse generator is a $\mathfrak g$-valued random variable $\mathring\omega^c$ with $\exp(\mathring\omega^c)=\mathring\Theta^c_e$ almost surely. For $\upsilon\in\mathfrak g$ put $\upsilon_{e'}=\mathring\Lambda_{e'}^{-1}\upsilon$ and $\upsilon_e=\mathring\Lambda_e^{-1}\upsilon$, let $\zeta_{e'}\sim\mathcal N(0,\mathring\Lambda_{e'}^{-1})$ and $\zeta_e\sim\mathcal N(0,\mathring\Lambda_e^{-1})$ be independent, and define the cubic and quartic commutator forms
\begin{equation}
\begin{aligned}
\kappa_3(\upsilon)={}&\langle\upsilon,[\upsilon_{e'},\upsilon_e]\rangle_{\mathsf G},\\
\kappa_4(\upsilon)={}&3\mathbb E\langle\upsilon,[\upsilon_{e'},\zeta_e]\rangle_{\mathsf G}^2+3\mathbb E\langle\upsilon,[\zeta_{e'},\upsilon_e]\rangle_{\mathsf G}^2\\
&+2\langle\upsilon,[\upsilon_{e'},[\upsilon_{e'},\upsilon_e]]+[\upsilon_e,[\upsilon_e,\upsilon_{e'}]]\rangle_{\mathsf G}.
\end{aligned}
\label{eq:llp_cumulant_forms}
\end{equation}

\begin{proposition}[Gaussian generator laws are not closed under coarse links]\label{prop:llp_coarse}
(i) If $\mathfrak g$ is abelian, $\mathring\omega_{e'}+\mathring\omega_e$ is a Gaussian coarse generator.
(ii) Let $G=\mathrm{GL}^+(K)$ with $K\geq2$, and let $\mathbb X_K$ be the set of real $K\times K$ matrices with $K$ distinct real eigenvalues, exactly two of them negative. Then $\mathbb X_K$ is a nonempty open subset of $\mathrm{GL}^+(K)$, no element of $\mathbb X_K$ has a real logarithm, and $\Pr(\mathring\Theta^c_e\in\mathbb X_K)>0$. Consequently no coarse generator exists, and the law of $\mathring\Theta^c_e$ is the image under the exponential of no law on $\mathfrak g$.
(iii) Let the laws depend on a width $\tau\in(0,1]$ through the means $\tau\mathring\mu_{e'},\tau\mathring\mu_e$ and the precisions $\tau^{-2}\mathring\Lambda_{e'},\tau^{-2}\mathring\Lambda_e$, with $\mathring\mu_{e'},\mathring\mu_e,\mathring\Lambda_{e'},\mathring\Lambda_e$ held. On the event $\mathsf E_\tau=\{\|\mathring\omega_{e'}\|+\|\mathring\omega_e\|\leq\frac12\}$ the principal logarithm of $\mathring\Theta^c_e$ exists and lies in $\mathfrak g$, and $\Pr(\mathsf E_\tau)\geq1-c_1e^{-c_2/\tau^2}$. If $\kappa_3$ is not identically zero, or if $\mathring\mu_{e'}=\mathring\mu_e=0$ and $\kappa_4$ is not identically zero, there is $\tau_0>0$ such that for every $\tau\in(0,\tau_0)$ no Gaussian random element of $\mathfrak{gl}(K,\mathbb R)$ coincides with $\operatorname{Log}\mathring\Theta^c_e$ almost surely on $\mathsf E_\tau$.
(iv) For $\mathfrak g=\mathfrak{gl}(K,\mathbb R)$ with $K\geq2$, $\mathsf G=I_K$ and $\mathring\Lambda_{e'}=\mathring\Lambda_e=\mathrm{Id}$, the form $\kappa_3$ vanishes identically and $\kappa_4(E_{12})=12$. With $\mathring\Lambda_{e'}=\mathrm{Id}$ and $\mathring\Lambda_e^{-1}$ the identity plus the orthogonal projection onto the line of $E_{12}$, $\kappa_3(E_{11}+E_{12})=1$. For an inner product on $\mathfrak g$ invariant under the adjoint action of $\mathfrak g$ and precisions that are multiples of the identity, both forms vanish identically.
\end{proposition}

\begin{proof}
(i) Commuting matrices satisfy $e^{\mathring\omega_{e'}}e^{\mathring\omega_e}=e^{\mathring\omega_{e'}+\mathring\omega_e}$, and a sum of independent Gaussian random elements is Gaussian.

(ii) The determinant of a matrix in $\mathbb X_K$ is the product of its eigenvalues, hence positive. Near such a matrix the eigenvalues depend continuously on the entries; a simple real eigenvalue stays real, because a nonreal eigenvalue nearby would bring its complex conjugate into the same small disk, and signs and distinctness persist, so $\mathbb X_K$ is open. Each negative eigenvalue is simple, so its single Jordan block occurs once, while Culver's criterion requires every Jordan block of a negative eigenvalue to occur an even number of times~\cite{culver1966existence}; hence no element of $\mathbb X_K$ has a real logarithm. With $\mathfrak J=E_{21}-E_{12}$, $\exp(\pi\mathfrak J)\exp(\operatorname{diag}(0,\log2,\log3,\ldots,\log K))=\operatorname{diag}(-1,-2,3,\ldots,K)\in\mathbb X_K$, so the preimage of $\mathbb X_K$ under the continuous map $(\xi_1,\xi_2)\mapsto e^{\xi_1}e^{\xi_2}$ is a nonempty open subset of $\mathfrak g\times\mathfrak g$, which the joint law of the generators charges because it has an everywhere positive density. A coarse generator would place $\mathring\Theta^c_e$ almost surely in $\exp(\mathfrak g)$, which misses $\mathbb X_K$, and the exponential image of any law on $\mathfrak g$ gives $\mathbb X_K$ mass zero.

(iii) Expansion of the logarithm. Let $\xi_1,\xi_2\in\mathfrak{gl}(K,\mathbb R)$ with $s=\|\xi_1\|+\|\xi_2\|\leq\frac12$, and let $\Pi_l=\sum_{a+b=l}\xi_1^a\xi_2^b/(a!b!)$ be the part of degree $l$ of $e^{\xi_1}e^{\xi_2}$. Then $\|e^{\xi_1}e^{\xi_2}-I\|\leq e^s-1\leq e^{1/2}-1<1$, so $\sum_{k\geq1}(-1)^{k+1}(e^{\xi_1}e^{\xi_2}-I)^k/k$ converges to the principal logarithm~\cite{Higham2008}. Substituting $\sum_{l\geq1}\Pi_l$ and grouping by total degree is legitimate, because replacing every coefficient by its absolute value and every word in $\xi_1,\xi_2$ by the product of the norms of its letters gives the expansion of $-\log(2-e^s)$, which converges since $\frac12<\log2$. Hence $\operatorname{Log}(e^{\xi_1}e^{\xi_2})=\sum_{l\geq1}\operatorname{BCH}_l(\xi_1,\xi_2)$ with homogeneous polynomials of degree $l$ satisfying $\|\operatorname{BCH}_l\|\leq b_ls^l$, where $b_l\geq0$ are the Taylor coefficients of $u\mapsto-\log(2-e^u)$ at zero. The first three are $\operatorname{BCH}_1=\Pi_1$, $\operatorname{BCH}_2=\Pi_2-\frac12\Pi_1^2=\frac12[\xi_1,\xi_2]$ and $\operatorname{BCH}_3=\Pi_3-\frac12(\Pi_1\Pi_2+\Pi_2\Pi_1)+\frac13\Pi_1^3=\frac1{12}([\xi_1,[\xi_1,\xi_2]]+[\xi_2,[\xi_2,\xi_1]])$, the Baker--Campbell--Hausdorff terms~\cite{Hall2015}, so
\begin{equation}
\begin{gathered}
\operatorname{Log}(e^{\xi_1}e^{\xi_2})=\xi_1+\xi_2+\tfrac12[\xi_1,\xi_2]+\tfrac1{12}\bigl([\xi_1,[\xi_1,\xi_2]]+[\xi_2,[\xi_2,\xi_1]]\bigr)+\operatorname{BCH}_{\geq4}(\xi_1,\xi_2),\\
\|\operatorname{BCH}_{\geq4}(\xi_1,\xi_2)\|\leq c_3s^4,\qquad c_3=\sum_{l\geq4}b_l2^{4-l}<\infty.
\end{gathered}
\label{eq:llp_bch}
\end{equation}

The logarithm lies in $\mathfrak g$ when $\xi_1,\xi_2\in\mathfrak g$. For $t\in[0,1]$ put $\gamma(t)=\operatorname{Log}(e^{t\xi_1}e^{\xi_2})$. By the bound above, $\gamma$ is real-analytic with $\|\gamma(t)\|\leq-\log(2-e^{1/2})<\frac\pi2$, and $\gamma(0)=\xi_2$, because the grouped series of $\operatorname{Log}(e^{\xi_2})$ is that of the formal identity $\log(e^u)=u$. Differentiating $e^{\gamma(t)}=e^{t\xi_1}e^{\xi_2}$ with Lemma~\ref{lem:llp_exp_derivative} gives $\mathbf L(\gamma,\gamma')=\xi_1e^\gamma$. Since $u\mapsto e^{u\gamma}\mathbf Te^{-u\gamma}$ and $u\mapsto e^{u\operatorname{ad}_\gamma}\mathbf T$, with $\operatorname{ad}_\gamma\mathbf T=[\gamma,\mathbf T]$, solve the same linear equation with the same initial value, $\mathbf L(\gamma,\mathbf T)=(\int_0^1e^{u\operatorname{ad}_\gamma}du)(\mathbf T)e^\gamma$. As $\|\operatorname{ad}_\gamma\|\leq2\|\gamma\|<\pi$, the operator $\int_0^1e^{u\operatorname{ad}_\gamma}du=\sum_{k\geq0}\operatorname{ad}_\gamma^k/(k+1)!$ is inverted by the power series of $w\mapsto w/(e^w-1)$, which converges for $|w|<2\pi$. Hence $\gamma$ solves the differential equation whose right side at $\mathbf T$ is that power series in $\operatorname{ad}_{\mathbf T}$ applied to $\xi_1$, an analytic vector field on the ball $\|\mathbf T\|<\pi$ that is tangent to $\mathfrak g$ at the points of $\mathfrak g$, because $\mathfrak g$ is a Lie subalgebra. The solution of this equation within $\mathfrak g$ from $\xi_2$ coincides with $\gamma$ by uniqueness while it exists, and it exists on $[0,1]$ because $\gamma$ stays in a compact subset of the ball; so $\operatorname{Log}(e^{\xi_1}e^{\xi_2})=\gamma(1)\in\mathfrak g$.

Probability of the event. Write $\mathring\omega_{e'}=\tau\chi_{e'}$ and $\mathring\omega_e=\tau\chi_e$ with $\chi_{e'}=\mathring\mu_{e'}+\zeta_{e'}$ and $\chi_e=\mathring\mu_e+\zeta_e$, independent with laws that do not depend on $\tau$. A Gaussian random element $\chi$ of a finite-dimensional space satisfies $\mathbb Ee^{c\|\chi\|^2}<\infty$ for some $c>0$, and Markov's inequality bounds $\Pr(\|\chi\|>1/(4\tau))$ by a multiple of $e^{-c/(16\tau^2)}$, which bounds the probability of the complement of $\mathsf E_\tau$.

Polynomial part. Put $\psi_1=\langle\upsilon,\chi_{e'}+\chi_e\rangle_{\mathsf G}$, $\psi_2=\frac12\langle\upsilon,[\chi_{e'},\chi_e]\rangle_{\mathsf G}$, $\psi_3=\frac1{12}\langle\upsilon,[\chi_{e'},[\chi_{e'},\chi_e]]+[\chi_e,[\chi_e,\chi_{e'}]]\rangle_{\mathsf G}$, $\psi_4=\tau^{-4}\langle\upsilon,\operatorname{BCH}_{\geq4}(\mathring\omega_{e'},\mathring\omega_e)\rangle_{\mathsf G}$ on $\mathsf E_\tau$ and $\psi_4=0$ off it, and $\Psi_\tau=\psi_1+\tau\psi_2+\tau^2\psi_3+\tau^3\psi_4$. By Eq.~\eqref{eq:llp_bch} and homogeneity, $\Psi_\tau=\tau^{-1}\langle\upsilon,\operatorname{Log}\mathring\Theta^c_e\rangle_{\mathsf G}$ on $\mathsf E_\tau$, and $|\psi_4|\leq c_4(\|\chi_{e'}\|+\|\chi_e\|)^4$ with $c_4$ independent of $\tau$, so every moment of $\psi_1,\ldots,\psi_4$ is bounded uniformly in $\tau$. Joint cumulants are multilinear and are polynomials in joint moments, and the cumulants of order three and four of the Gaussian variable $\psi_1$ vanish, so, with $\operatorname{cum}_l$ the cumulant of order $l$ and $\operatorname{cum}$ the joint cumulant,
\begin{equation}
\begin{aligned}
\operatorname{cum}_3\Psi_\tau={}&3\tau\operatorname{cum}(\psi_1,\psi_1,\psi_2)+O(\tau^2),\\
\operatorname{cum}_4\Psi_\tau={}&4\tau\operatorname{cum}(\psi_1,\psi_1,\psi_1,\psi_2)\\
&+\tau^2\bigl[6\operatorname{cum}(\psi_1,\psi_1,\psi_2,\psi_2)+4\operatorname{cum}(\psi_1,\psi_1,\psi_1,\psi_3)\bigr]+O(\tau^3).
\end{aligned}
\label{eq:llp_cumulant_expansion}
\end{equation}

Gaussian integration by parts. Let $\partial_\upsilon$ be the derivative of a function of $(\zeta_{e'},\zeta_e)$ along $(\upsilon_{e'},\upsilon_e)$. Every polynomial $\varphi$ of $(\zeta_{e'},\zeta_e)$ satisfies $\mathbb E[(\psi_1-\mathbb E\psi_1)\varphi]=\mathbb E[\partial_\upsilon\varphi]$, by Eq.~\eqref{eq:llp_stein} applied to the two summands of $\psi_1-\mathbb E\psi_1=\langle\upsilon,\zeta_{e'}\rangle_{\mathsf G}+\langle\upsilon,\zeta_e\rangle_{\mathsf G}$. Two applications give $\operatorname{cum}(\psi_1,\psi_1,\psi_2)=\mathbb E[(\psi_1-\mathbb E\psi_1)^2(\psi_2-\mathbb E\psi_2)]=\partial_\upsilon^2\psi_2=\kappa_3(\upsilon)$ for all means, because $\psi_2$ is a quadratic polynomial whose second derivative along $(\upsilon_{e'},\upsilon_e)$ is $\langle\upsilon,[\upsilon_{e'},\upsilon_e]\rangle_{\mathsf G}$. If the means vanish, the map $(\zeta_{e'},\zeta_e)\mapsto(-\zeta_{e'},-\zeta_e)$ preserves the law and changes the signs of $\psi_1$ and $\psi_3$ but not of $\psi_2$, so every joint moment and cumulant with an odd number of entries from $\{\psi_1,\psi_3\}$ vanishes; in particular $\operatorname{cum}(\psi_1,\psi_1,\psi_1,\psi_2)=0$, while $\mathbb E\psi_2=0$ because the generators are independent and centered. Repeated integration by parts gives $\mathbb E[\psi_1^2\psi_2^2]=\mathbb E[\psi_1^2]\mathbb E[\psi_2^2]+2\mathbb E[(\partial_\upsilon\psi_2)^2]$ and $\mathbb E[\psi_1^3\psi_3]=3\mathbb E[\psi_1^2]\mathbb E[\psi_1\psi_3]+\partial_\upsilon^3\psi_3$, hence $\operatorname{cum}(\psi_1,\psi_1,\psi_2,\psi_2)=2\mathbb E[(\partial_\upsilon\psi_2)^2]$ and $\operatorname{cum}(\psi_1,\psi_1,\psi_1,\psi_3)=\partial_\upsilon^3\psi_3$, which is six times the value of the homogeneous cubic $\psi_3$ at $(\upsilon_{e'},\upsilon_e)$. Since $\partial_\upsilon\psi_2=\frac12\langle\upsilon,[\upsilon_{e'},\zeta_e]\rangle_{\mathsf G}+\frac12\langle\upsilon,[\zeta_{e'},\upsilon_e]\rangle_{\mathsf G}$ has independent centered summands, the coefficient of $\tau^2$ in Eq.~\eqref{eq:llp_cumulant_expansion} is $\kappa_4(\upsilon)$ in the centered case.

Gaussian rigidity. Fix $\upsilon$ with $\kappa_3(\upsilon)\neq0$, or, in the centered case, with $\kappa_4(\upsilon)\neq0$, and suppose that along a sequence of widths tending to zero some Gaussian random element of $\mathfrak{gl}(K,\mathbb R)$ coincides with $\operatorname{Log}\mathring\Theta^c_e$ almost surely on $\mathsf E_\tau$. Pairing it with $\upsilon$ and dividing by $\tau$ gives real Gaussian variables $\mathrm X_\tau\sim\mathcal N(m_\tau,s_\tau^2)$ equal to $\Psi_\tau$ outside an event of probability at most $c_1e^{-c_2/\tau^2}$. Since $\Psi_\tau\to\psi_1$ in mean square, $\mathrm X_\tau\to\psi_1$ in probability, hence in law. The moduli $e^{-s_\tau^2t^2/2}$ of the characteristic functions then converge, so $s_\tau^2\to\operatorname{Var}\psi_1$, and the phases $e^{im_\tau t}$ converge to $e^{it\mathbb E\psi_1}$ for every $t$; integrating $e^{i(m_\tau-\mathbb E\psi_1)t}$ against the standard normal density gives $e^{-(m_\tau-\mathbb E\psi_1)^2/2}\to1$, so $m_\tau\to\mathbb E\psi_1$. The moments of order eight of $\mathrm X_\tau$ are therefore bounded, and the Cauchy--Schwarz inequality on the exceptional event shows that the moments of order at most four of $\mathrm X_\tau$ and $\Psi_\tau$ differ by $O(e^{-c_2/(2\tau^2)})$, and so do their cumulants of order three and four. Those of $\mathrm X_\tau$ vanish, so by Eq.~\eqref{eq:llp_cumulant_expansion} the quantity $3\tau\kappa_3(\upsilon)+O(\tau^2)$, and in the centered case $\tau^2\kappa_4(\upsilon)+O(\tau^3)$, would tend to zero faster than every power of $\tau$, which is impossible. Hence such a sequence does not exist, and $\tau_0$ does.

(iv) For $\mathsf G=I_K$ the inner product is $\operatorname{tr}(\mathbf S^\top\mathbf T)$. With unit precisions $\upsilon_{e'}=\upsilon_e=\upsilon$, so $\kappa_3(\upsilon)=\langle\upsilon,[\upsilon,\upsilon]\rangle=0$ and the double brackets of $\kappa_4$ vanish. For $\upsilon=E_{12}$, $\langle\upsilon,[\upsilon,\zeta]\rangle=\zeta_{22}-\zeta_{11}$, of variance $2$ when $\zeta$ has independent standard normal entries, and $\langle\upsilon,[\zeta,\upsilon]\rangle$ is its negative, so $\kappa_4(\upsilon)=3\cdot2+3\cdot2=12$. In the second instance $\upsilon_{e'}=\upsilon=E_{11}+E_{12}$ and $\upsilon_e=E_{11}+2E_{12}$, whose bracket $[\upsilon_{e'},\upsilon_e]=E_{12}$ pairs with $\upsilon$ to $1$. For an invariant inner product, $\langle\upsilon,[\upsilon,\mathbf T]\rangle_{\mathsf G}=-\langle[\upsilon,\upsilon],\mathbf T\rangle_{\mathsf G}=0$ for every $\mathbf T\in\mathfrak g$, and with precisions that are multiples of the identity $\upsilon_{e'}$ and $\upsilon_e$ are multiples of $\upsilon$, so every term of both forms vanishes.
\end{proof}

When the logarithm exists, the commutator term of the expansion thus makes the coarse generator non-Gaussian at first order in the width whenever $\kappa_3$ is not identically zero and, for centered laws, at second order whenever $\kappa_4$ is not. For an invariant inner product with isotropic precisions both leading coefficients vanish, and this test gives no conclusion.

\begin{remark}[A finite-dimensional description of coarse link laws]\label{rem:llp_open_coarse}
The law of the coarse link always exists, as the pushforward of the fine law through the ordered product~\eqref{eq:pifb_coarse_link_product}, and the exact coarse-graining of Section~\ref{sec:rg_status} applies to it. For independent links in series the step sends laws $\nu'$ and $\nu$ on $G$ to the law $\nu'*\nu$ of their product, so a finite-dimensional family of link laws that survives the step must be closed under this convolution. For $G=\mathrm{GL}^+(K)$ with $K\geq2$ the question is open. It would be settled by a weakly continuous injective map from a connected open subset of some $\mathbb R^d$ into the Borel probability laws on $\mathrm{GL}^+(K)$ whose image contains the exponential image of every nondegenerate Gaussian law on $\mathfrak{gl}(K,\mathbb R)$ and contains $\nu'*\nu$ whenever it contains $\nu'$ and $\nu$, or by a proof that for some $K\geq2$ no such map exists.
\end{remark}
